\documentclass[8pt]{beamer}
\usetheme{Madrid}
\usecolortheme{default}
\usepackage{amsmath, amssymb, graphicx, array, multirow, booktabs, xcolor, tikz}
\usetikzlibrary{shapes,arrows,positioning,fit,backgrounds}

% Compact font settings
\setbeamerfont{title}{size=\Large}
\setbeamerfont{subtitle}{size=\small}
\setbeamerfont{author}{size=\scriptsize}
\setbeamerfont{date}{size=\scriptsize}
\setbeamerfont{frametitle}{size=\large}
\setbeamerfont{framesubtitle}{size=\normalsize}
\setbeamerfont{enumerate item}{size=\footnotesize}
\setbeamerfont{itemize item}{size=\footnotesize}
\setbeamerfont{block title}{size=\footnotesize}

% Compact spacing
\setlength{\itemsep}{0.08ex}
\setlength{\parskip}{0.08ex}
\setlength{\leftmargini}{0.3cm}
\setbeamersize{text margin left=3mm, text margin right=3mm}
\addtobeamertemplate{frame begin}{\vspace{-0.4cm}}{}

% Colors
\definecolor{myblue}{RGB}{0,0,255}
\definecolor{myred}{RGB}{255,0,0}
\definecolor{mygreen}{RGB}{0,128,0}
\definecolor{myorange}{RGB}{255,165,0}
\definecolor{mypurple}{RGB}{128,0,128}
\definecolor{mygray}{RGB}{128,128,128}
\definecolor{mygold}{RGB}{255,215,0}

\title{Generative Artificial Intelligence}
\subtitle{100 Questions: GenAI, LLMs, Transformers, Word2Vec, RNNs, LSTM}
\author{GARP RAI Exam Preparation}
\date{}

\begin{document}
	
	\begin{frame}
		\titlepage
	\end{frame}
	
	% ============================================================
	% SECTION 1: INTRODUCTION TO GENAI - QUESTIONS 1-10
	% ============================================================
	
	\begin{frame}{Q1: Generative AI Definition - Tutorial}
		\fontsize{6.5}{7.5}\selectfont
		\textbf{QUESTION: What is Generative Artificial Intelligence (GenAI)?}
		
		\vspace{0.1cm}
		\textbf{OPTIONS:}
		\begin{enumerate}
			\item A branch of AI that analyzes existing data to make predictions
			\item A branch of AI that employs models to generate text, sound, or videos
			\item A branch of AI that only classifies images
			\item A branch of AI that performs regression analysis
		\end{enumerate}
		
		\vspace{0.1cm}
		\textcolor{myblue}{\textbf{ANSWER: B}}
		
		\vspace{0.1cm}
		\textbf{DETAILED EXPLANATION:}
		\begin{itemize}
			\item \textbf{Definition:} GenAI creates new content rather than just analyzing existing data
			\item \textbf{Key characteristic:} Generates text, sound, or videos
			\item \textbf{Example:} ChatGPT generates text, DALL-E generates images
		\end{itemize}
		
		\vspace{0.1cm}
		\textcolor{myred}{\textbf{WHY OTHER OPTIONS ARE WRONG:}}
		\begin{itemize}
			\item \textcolor{myred}{A:} This describes traditional AI/machine learning - it analyzes, not generates
			\item \textcolor{myred}{C:} This is too narrow - GenAI creates more than just images
			\item \textcolor{myred}{D:} Regression is a traditional ML technique, not GenAI
		\end{itemize}
		
		\textcolor{myred}{\textbf{EXAM TRAP:}} GenAI is about \textcolor{myblue}{creating} new content, not just analyzing!
	\end{frame}
	
	\begin{frame}{Q2: ChatGPT Milestone - Tutorial}
		\fontsize{6.5}{7.5}\selectfont
		\textbf{QUESTION: What was significant about ChatGPT's release in November 2022?}
		
		\vspace{0.1cm}
		\textbf{OPTIONS:}
		\begin{enumerate}
			\item It was the first AI tool ever created
			\item It reached 1 million users in just 5 days, making it the fastest AI tool to do so
			\item It was the first chatbot
			\item It had no impact on AI adoption
		\end{enumerate}
		
		\vspace{0.1cm}
		\textcolor{myblue}{\textbf{ANSWER: B}}
		
		\vspace{0.1cm}
		\textbf{DETAILED EXPLANATION:}
		\begin{itemize}
			\item ChatGPT reached 1 million users in just 5 days
			\item Facebook took almost a year to reach this milestone
			\item Demonstrated rapid adoption of GenAI
			\item Sparked widespread interest and investment
		\end{itemize}
		
		\vspace{0.1cm}
		\textcolor{myred}{\textbf{WHY OTHER OPTIONS ARE WRONG:}}
		\begin{itemize}
			\item \textcolor{myred}{A:} AI tools existed long before ChatGPT
			\item \textcolor{myred}{C:} Chatbots existed before ChatGPT (ELIZA in 1960s)
			\item \textcolor{myred}{D:} ChatGPT had enormous impact on AI adoption
		\end{itemize}
		
		\textcolor{myred}{\textbf{EXAM TRAP:}} ChatGPT was the \textcolor{myblue}{fastest AI tool} to reach 1 million users!
	\end{frame}
	
	\begin{frame}{Q3: GenAI vs Traditional AI - Tutorial}
		\fontsize{6.5}{7.5}\selectfont
		\textbf{QUESTION: What is the key difference between Generative AI and Traditional AI?}
		
		\vspace{0.1cm}
		\textbf{OPTIONS:}
		\begin{enumerate}
			\item Traditional AI creates new content; GenAI analyzes existing data
			\item GenAI creates new content; Traditional AI analyzes existing data and makes predictions
			\item Both create new content equally
			\item Neither creates new content
		\end{enumerate}
		
		\vspace{0.1cm}
		\textcolor{myblue}{\textbf{ANSWER: B}}
		
		\vspace{0.1cm}
		\textbf{DETAILED EXPLANATION:}
		\begin{itemize}
			\item \textbf{Traditional AI:} Analyzes existing data, makes predictions/classifications
			\item \textbf{GenAI:} Creates new content (text, images, audio, video)
			\item \textbf{Traditional AI outputs:} Labels or numerical values
			\item \textbf{GenAI outputs:} Novel, creative content
		\end{itemize}
		
		\vspace{0.1cm}
		\textcolor{myred}{\textbf{WHY OTHER OPTIONS ARE WRONG:}}
		\begin{itemize}
			\item \textcolor{myred}{A:} This reverses the definitions
			\item \textcolor{myred}{C:} They are fundamentally different in output
			\item \textcolor{myred}{D:} GenAI definitely creates new content
		\end{itemize}
		
		\textcolor{myred}{\textbf{EXAM TRAP:}} GenAI = \textcolor{myblue}{creates}; Traditional AI = \textcolor{myblue}{analyzes}!
	\end{frame}
	
	\begin{frame}{Q4: GenAI Content Modalities - Tutorial}
		\fontsize{6.5}{7.5}\selectfont
		\textbf{QUESTION: What are the five major content modalities in GenAI?}
		
		\vspace{0.1cm}
		\textbf{OPTIONS:}
		\begin{enumerate}
			\item Text, Image, Audio, Video, and Multimodal
			\item Text, Image, Sound, Animation, and 3D
			\item Numbers, Text, Images, Audio, and Video
			\item Text, Photos, Music, Movies, and Art
		\end{enumerate}
		
		\vspace{0.1cm}
		\textcolor{myblue}{\textbf{ANSWER: A}}
		
		\vspace{0.1cm}
		\textbf{DETAILED EXPLANATION:}
		\begin{itemize}
			\item \textbf{Text:} ChatGPT, Claude, Gemini
			\item \textbf{Image:} DALL-E, Stable Diffusion
			\item \textbf{Audio:} Text-to-speech, music generation
			\item \textbf{Video:} Sora, RunwayML
			\item \textbf{Multimodal:} Gemini, GPT-4 (multiple input/output types)
		\end{itemize}
		
		\vspace{0.1cm}
		\textcolor{myred}{\textbf{WHY OTHER OPTIONS ARE WRONG:}}
		\begin{itemize}
			\item \textcolor{myred}{B:} "Animation" and "3D" are subcategories, not major modalities
			\item \textcolor{myred}{C:} "Numbers" is traditional data, not a GenAI modality
			\item \textcolor{myred}{D:} Uses casual terms, not the technical classification
		\end{itemize}
		
		\textcolor{myred}{\textbf{EXAM TRAP:}} The five modalities are \textcolor{myblue}{Text, Image, Audio, Video, Multimodal}!
	\end{frame}
	
	\begin{frame}{Q5: Multimodal Models - Tutorial}
		\fontsize{6.5}{7.5}\selectfont
		\textbf{QUESTION: What does a multimodal GenAI model do?}
		
		\vspace{0.1cm}
		\textbf{OPTIONS:}
		\begin{enumerate}
			\item Only processes text inputs
			\item Takes multiple input types and generates multiple output types
			\item Only generates images
			\item Only processes a single type of data
		\end{enumerate}
		
		\vspace{0.1cm}
		\textcolor{myblue}{\textbf{ANSWER: B}}
		
		\vspace{0.1cm}
		\textbf{DETAILED EXPLANATION:}
		\begin{itemize}
			\item Multimodal models accept various input types
			\item Can generate output in different formats
			\item Example: Gemini can take text and images as input
			\item Example: GPT-4 can analyze images and generate text
		\end{itemize}
		
		\vspace{0.1cm}
		\textcolor{myred}{\textbf{WHY OTHER OPTIONS ARE WRONG:}}
		\begin{itemize}
			\item \textcolor{myred}{A:} This describes text-only models (unimodal)
			\item \textcolor{myred}{C:} This describes image-only models
			\item \textcolor{myred}{D:} This describes unimodal models, not multimodal
		\end{itemize}
		
		\textcolor{myred}{\textbf{EXAM TRAP:}} Multimodal = \textcolor{myblue}{multiple input and output types}!
	\end{frame}
	
	\begin{frame}{Q6: LLM Definition - Tutorial}
		\fontsize{6.5}{7.5}\selectfont
		\textbf{QUESTION: What is a Large Language Model (LLM)?}
		
		\vspace{0.1cm}
		\textbf{OPTIONS:}
		\begin{enumerate}
			\item A small neural network with few parameters
			\item A large neural network with billions of parameters, trained on vast text data
			\item A model that only generates images
			\item A model that only processes numerical data
		\end{enumerate}
		
		\vspace{0.1cm}
		\textcolor{myblue}{\textbf{ANSWER: B}}
		
		\vspace{0.1cm}
		\textbf{DETAILED EXPLANATION:}
		\begin{itemize}
			\item LLMs have billions to trillions of parameters
			\item Trained on very large text corpora
			\item Use transformer architecture
			\item Examples: GPT-3 (175B parameters), GPT-4
		\end{itemize}
		
		\vspace{0.1cm}
		\textcolor{myred}{\textbf{WHY OTHER OPTIONS ARE WRONG:}}
		\begin{itemize}
			\item \textcolor{myred}{A:} LLMs are LARGE, not small
			\item \textcolor{myred}{C:} LLMs are text-focused, not image generation
			\item \textcolor{myred}{D:} LLMs process text, not primarily numerical data
		\end{itemize}
		
		\textcolor{myred}{\textbf{EXAM TRAP:}} LLMs are \textcolor{myblue}{large neural networks} with billions of parameters!
	\end{frame}
	
	\begin{frame}{Q7: Transformers Introduction - Tutorial}
		\fontsize{6.5}{7.5}\selectfont
		\textbf{QUESTION: What is the key innovation of transformers in NLP?}
		
		\vspace{0.1cm}
		\textbf{OPTIONS:}
		\begin{enumerate}
			\item They use only feedforward networks
			\item They use the "attention" mechanism to understand word relationships
			\item They are slower than RNNs
			\item They process text sequentially
		\end{enumerate}
		
		\vspace{0.1cm}
		\textcolor{myblue}{\textbf{ANSWER: B}}
		
		\vspace{0.1cm}
		\textbf{DETAILED EXPLANATION:}
		\begin{itemize}
			\item Transformers use self-attention mechanism
			\item Attention allows understanding relationships between all words
			\item Process entire sequence in parallel
			\item Foundation for modern LLMs
		\end{itemize}
		
		\vspace{0.1cm}
		\textcolor{myred}{\textbf{WHY OTHER OPTIONS ARE WRONG:}}
		\begin{itemize}
			\item \textcolor{myred}{A:} Transformers use attention AND feedforward networks
			\item \textcolor{myred}{C:} Transformers are faster than RNNs (parallel processing)
			\item \textcolor{myred}{D:} Transformers process in parallel, not sequentially
		\end{itemize}
		
		\textcolor{myred}{\textbf{EXAM TRAP:}} Transformers = \textcolor{myblue}{attention mechanism}!
	\end{frame}
	
	\begin{frame}{Q8: DALL-E Function - Tutorial}
		\fontsize{6.5}{7.5}\selectfont
		\textbf{QUESTION: What does DALL-E do?}
		
		\vspace{0.1cm}
		\textbf{OPTIONS:}
		\begin{enumerate}
			\item Generates text from images
			\item Creates images from text prompts
			\item Generates audio from text
			\item Analyzes financial data
		\end{enumerate}
		
		\vspace{0.1cm}
		\textcolor{myblue}{\textbf{ANSWER: B}}
		
		\vspace{0.1cm}
		\textbf{DETAILED EXPLANATION:}
		\begin{itemize}
			\item DALL-E is developed by OpenAI
			\item Takes text prompts and creates images
			\item Can generate novel, creative images
			\item Powerful image generation tool
		\end{itemize}
		
		\vspace{0.1cm}
		\textcolor{myred}{\textbf{WHY OTHER OPTIONS ARE WRONG:}}
		\begin{itemize}
			\item \textcolor{myred}{A:} This describes image-to-text (opposite of DALL-E)
			\item \textcolor{myred}{C:} Audio generation is a different modality
			\item \textcolor{myred}{D:} Financial analysis is not DALL-E's function
		\end{itemize}
		
		\textcolor{myred}{\textbf{EXAM TRAP:}} DALL-E = \textcolor{myblue}{image generation from text prompts}!
	\end{frame}
	
	\begin{frame}{Q9: GANs Architecture - Tutorial}
		\fontsize{6.5}{7.5}\selectfont
		\textbf{QUESTION: What are the two competing networks in a Generative Adversarial Network (GAN)?}
		
		\vspace{0.1cm}
		\textbf{OPTIONS:}
		\begin{enumerate}
			\item Encoder and Decoder
			\item Generator and Discriminator
			\item Input and Output
			\item Training and Testing
		\end{enumerate}
		
		\vspace{0.1cm}
		\textcolor{myblue}{\textbf{ANSWER: B}}
		
		\vspace{0.1cm}
		\textbf{DETAILED EXPLANATION:}
		\begin{itemize}
			\item \textbf{Generator:} Creates simulated data, learns to fool discriminator
			\item \textbf{Discriminator:} Identifies real vs fake data, gives feedback
			\item Process iterates until generated data is indistinguishable
			\item Used for images, audio, video creation
		\end{itemize}
		
		\vspace{0.1cm}
		\textcolor{myred}{\textbf{WHY OTHER OPTIONS ARE WRONG:}}
		\begin{itemize}
			\item \textcolor{myred}{A:} Encoder/Decoder are in VAEs, not GANs
			\item \textcolor{myred}{C:} Input/Output are general terms, not GAN-specific
			\item \textcolor{myred}{D:} Training/Testing are ML phases, not GAN architecture
		\end{itemize}
		
		\textcolor{myred}{\textbf{EXAM TRAP:}} GANs = \textcolor{myblue}{Generator + Discriminator} competing!
	\end{frame}
	
	\begin{frame}{Q10: GenAI vs LLM Relationship - Tutorial}
		\fontsize{6.5}{7.5}\selectfont
		\textbf{QUESTION: Which statement correctly describes the relationship between GenAI and LLMs?}
		
		\vspace{0.1cm}
		\textbf{OPTIONS:}
		\begin{enumerate}
			\item GenAI and LLMs are the same thing
			\item LLMs are a subset of GenAI that focus on text generation
			\item GenAI is a subset of LLMs
			\item They are completely unrelated
		\end{enumerate}
		
		\vspace{0.1cm}
		\textcolor{myblue}{\textbf{ANSWER: B}}
		
		\vspace{0.1cm}
		\textbf{DETAILED EXPLANATION:}
		\begin{itemize}
			\item GenAI is broader - includes images, audio, video
			\item LLMs specifically focus on text generation
			\item All LLMs are GenAI, but not all GenAI are LLMs
			\item GenAI = umbrella term; LLMs = specific type
		\end{itemize}
		
		\vspace{0.1cm}
		\textcolor{myred}{\textbf{WHY OTHER OPTIONS ARE WRONG:}}
		\begin{itemize}
			\item \textcolor{myred}{A:} GenAI is broader than just LLMs
			\item \textcolor{myred}{C:} This reverses the relationship
			\item \textcolor{myred}{D:} They are related - LLMs are a type of GenAI
		\end{itemize}
		
		\textcolor{myred}{\textbf{EXAM TRAP:}} All LLMs are GenAI, but \textcolor{myblue}{not all GenAI are LLMs}!
	\end{frame}
	
	% ============================================================
	% SECTION 2: WORD EMBEDDINGS & WORD2VEC - QUESTIONS 11-25
	% ============================================================
	
	\begin{frame}{Q11: BoW Limitation - Tutorial}
		\fontsize{6.5}{7.5}\selectfont
		\textbf{QUESTION: What is a major limitation of the Bag of Words (BoW) approach?}
		
		\vspace{0.1cm}
		\textbf{OPTIONS:}
		\begin{enumerate}
			\item It captures semantic meaning of words
			\item It creates dense vectors with few zeros
			\item It treats words independently and loses context
			\item It handles negation words perfectly
		\end{enumerate}
		
		\vspace{0.1cm}
		\textcolor{myblue}{\textbf{ANSWER: C}}
		
		\vspace{0.1cm}
		\textbf{DETAILED EXPLANATION:}
		\begin{itemize}
			\item BoW ignores word order and context
			\item "not good" and "good" both contain "good"
			\item Cannot identify similar meanings ("amazing" vs "fantastic")
			\item Creates sparse vectors with many zeros
		\end{itemize}
		
		\vspace{0.1cm}
		\textcolor{myred}{\textbf{WHY OTHER OPTIONS ARE WRONG:}}
		\begin{itemize}
			\item \textcolor{myred}{A:} BoW does NOT capture semantic meaning
			\item \textcolor{myred}{B:} BoW creates SPARSE, not dense vectors
			\item \textcolor{myred}{D:} BoW struggles with negation words
		\end{itemize}
		
		\textcolor{myred}{\textbf{EXAM TRAP:}} BoW loses \textcolor{myblue}{context and semantic meaning}!
	\end{frame}
	
	\begin{frame}{Q12: Word2Vec Creator - Tutorial}
		\fontsize{6.5}{7.5}\selectfont
		\textbf{QUESTION: Who created Word2Vec and when?}
		
		\vspace{0.1cm}
		\textbf{OPTIONS:}
		\begin{enumerate}
			\item OpenAI in 2020
			\item Google team led by Tomas Mikolov in 2013
			\item Facebook in 2018
			\item Microsoft in 2015
		\end{enumerate}
		
		\vspace{0.1cm}
		\textcolor{myblue}{\textbf{ANSWER: B}}
		
		\vspace{0.1cm}
		\textbf{DETAILED EXPLANATION:}
		\begin{itemize}
			\item Word2Vec created by Google team (2013)
			\item Led by Tomas Mikolov
			\item Alternative to conventional BoW
			\item Revolutionary for NLP
		\end{itemize}
		
		\vspace{0.1cm}
		\textcolor{myred}{\textbf{WHY OTHER OPTIONS ARE WRONG:}}
		\begin{itemize}
			\item \textcolor{myred}{A:} OpenAI was not involved in Word2Vec
			\item \textcolor{myred}{C:} Facebook created other models (BART, RoBERTa)
			\item \textcolor{myred}{D:} Microsoft was not involved in Word2Vec
		\end{itemize}
		
		\textcolor{myred}{\textbf{EXAM TRAP:}} Word2Vec = \textcolor{myblue}{Google 2013} by Tomas Mikolov!
	\end{frame}
	
	\begin{frame}{Q13: Word Embeddings Definition - Tutorial}
		\fontsize{6.5}{7.5}\selectfont
		\textbf{QUESTION: What are word embeddings in NLP?}
		
		\vspace{0.1cm}
		\textbf{OPTIONS:}
		\begin{enumerate}
			\item Sparse vectors with mostly zeros
			\item Dense vector representations of words that encode semantic information
			\item A list of synonyms for each word
			\item A method for removing stop words
		\end{enumerate}
		
		\vspace{0.1cm}
		\textcolor{myblue}{\textbf{ANSWER: B}}
		
		\vspace{0.1cm}
		\textbf{DETAILED EXPLANATION:}
		\begin{itemize}
			\item Dense vectors (fewer dimensions)
			\item Similar words have similar vectors
			\item Encode semantic meaning
			\item Generated through neural network training
		\end{itemize}
		
		\vspace{0.1cm}
		\textcolor{myred}{\textbf{WHY OTHER OPTIONS ARE WRONG:}}
		\begin{itemize}
			\item \textcolor{myred}{A:} Embeddings are DENSE, not sparse
			\item \textcolor{myred}{C:} Embeddings are vectors, not word lists
			\item \textcolor{myred}{D:} Embeddings are for representation, not removal
		\end{itemize}
		
		\textcolor{myred}{\textbf{EXAM TRAP:}} Embeddings = \textcolor{myblue}{dense vectors capturing meaning}!
	\end{frame}
	
	\begin{frame}{Q14: CBoW Architecture - Tutorial}
		\fontsize{6.5}{7.5}\selectfont
		\textbf{QUESTION: What is the Continuous Bag of Words (CBoW) architecture in Word2Vec?}
		
		\vspace{0.1cm}
		\textbf{OPTIONS:}
		\begin{enumerate}
			\item Predicts context words from a target word
			\item Predicts a target word from surrounding context words
			\item Predicts the next word in a sequence
			\item Only analyzes single words
		\end{enumerate}
		
		\vspace{0.1cm}
		\textcolor{myblue}{\textbf{ANSWER: B}}
		
		\vspace{0.1cm}
		\textbf{DETAILED EXPLANATION:}
		\begin{itemize}
			\item "Fill in the blank" technique
			\item Input: context words (before and after)
			\item Output: target word (the missing word)
			\item Goal: $P(w_i | w_{i-c}, ..., w_{i+c})$
		\end{itemize}
		
		\vspace{0.1cm}
		\textcolor{myred}{\textbf{WHY OTHER OPTIONS ARE WRONG:}}
		\begin{itemize}
			\item \textcolor{myred}{A:} This describes Skip-gram (context from target)
			\item \textcolor{myred}{C:} This describes autoregressive models
			\item \textcolor{myred}{D:} CBoW uses context, not just single words
		\end{itemize}
		
		\textcolor{myred}{\textbf{EXAM TRAP:}} CBoW = \textcolor{myblue}{context $\to$ target} word!
	\end{frame}
	
	\begin{frame}{Q15: Skip-gram Architecture - Tutorial}
		\fontsize{6.5}{7.5}\selectfont
		\textbf{QUESTION: What is the Skip-gram architecture in Word2Vec?}
		
		\vspace{0.1cm}
		\textbf{OPTIONS:}
		\begin{enumerate}
			\item Predicts a target word from surrounding context words
			\item Predicts surrounding context words from a target word
			\item Predicts the next sentence
			\item Analyzes word frequency only
		\end{enumerate}
		
		\vspace{0.1cm}
		\textcolor{myblue}{\textbf{ANSWER: B}}
		
		\vspace{0.1cm}
		\textbf{DETAILED EXPLANATION:}
		\begin{itemize}
			\item "Word association" technique
			\item Input: target word
			\item Output: context words (surrounding words)
			\item Goal: $P(w_{i-c}, ..., w_{i+c} | w_i)$
		\end{itemize}
		
		\vspace{0.1cm}
		\textcolor{myred}{\textbf{WHY OTHER OPTIONS ARE WRONG:}}
		\begin{itemize}
			\item \textcolor{myred}{A:} This describes CBoW
			\item \textcolor{myred}{C:} Skip-gram predicts words, not sentences
			\item \textcolor{myred}{D:} Skip-gram uses neural networks, not just frequency
		\end{itemize}
		
		\textcolor{myred}{\textbf{EXAM TRAP:}} Skip-gram = \textcolor{myblue}{target $\to$ context} words!
	\end{frame}
	
	\begin{frame}{Q16: CBoW vs Skip-gram - Tutorial}
		\fontsize{6.5}{7.5}\selectfont
		\textbf{QUESTION: What is the key difference between CBoW and Skip-gram?}
		
		\vspace{0.1cm}
		\textbf{OPTIONS:}
		\begin{enumerate}
			\item CBoW is for English; Skip-gram for other languages
			\item CBoW predicts target from context; Skip-gram predicts context from target
			\item CBoW is slower than Skip-gram
			\item There is no difference - they are the same
		\end{enumerate}
		
		\vspace{0.1cm}
		\textcolor{myblue}{\textbf{ANSWER: B}}
		
		\vspace{0.1cm}
		\textbf{DETAILED EXPLANATION:}
		\begin{itemize}
			\item \textbf{CBoW:} context $\to$ target
			\item \textbf{Skip-gram:} target $\to$ context
			\item Both are Word2Vec architectures
			\item Skip-gram better for rare words; CBoW faster
		\end{itemize}
		
		\vspace{0.1cm}
		\textcolor{myred}{\textbf{WHY OTHER OPTIONS ARE WRONG:}}
		\begin{itemize}
			\item \textcolor{myred}{A:} Both work for all languages
			\item \textcolor{myred}{C:} CBoW is actually faster
			\item \textcolor{myred}{D:} They are distinctly different architectures
		\end{itemize}
		
		\textcolor{myred}{\textbf{EXAM TRAP:}} CBoW = context $\to$ target; Skip-gram = \textcolor{myblue}{target $\to$ context}!
	\end{frame}
	
	\begin{frame}{Q17: Word2Vec Dimensionality - Tutorial}
		\fontsize{6.5}{7.5}\selectfont
		\textbf{QUESTION: How does Word2Vec reduce dimensionality compared to BoW?}
		
		\vspace{0.1cm}
		\textbf{OPTIONS:}
		\begin{enumerate}
			\item Word2Vec uses more dimensions than BoW
			\item Word2Vec uses dense 300-dimensional vectors instead of vocabulary-sized sparse vectors
			\item Word2Vec doesn't reduce dimensionality
			\item Word2Vec uses one-hot encoding
		\end{enumerate}
		
		\vspace{0.1cm}
		\textcolor{myblue}{\textbf{ANSWER: B}}
		
		\vspace{0.1cm}
		\textbf{DETAILED EXPLANATION:}
		\begin{itemize}
			\item Traditional BoW: vocabulary size dimensions (e.g., 100,000)
			\item Word2Vec: typically 300 dimensions
			\item Reduction: 99.7\% fewer weights!
			\item Dense vs sparse representation
		\end{itemize}
		
		\vspace{0.1cm}
		\textcolor{myred}{\textbf{WHY OTHER OPTIONS ARE WRONG:}}
		\begin{itemize}
			\item \textcolor{myred}{A:} Word2Vec uses FEWER dimensions
			\item \textcolor{myred}{C:} Dimensionality reduction is a key feature
			\item \textcolor{myred}{D:} Word2Vec uses embeddings, not one-hot encoding
		\end{itemize}
		
		\textcolor{myred}{\textbf{EXAM TRAP:}} Word2Vec reduces dimensions from \textcolor{myblue}{vocabulary size to ~300}!
	\end{frame}
	
	\begin{frame}{Q18: Cosine Similarity Formula - Tutorial}
		\fontsize{6.5}{7.5}\selectfont
		\textbf{QUESTION: What is the formula for cosine similarity between two vectors P and Q?}
		
		\vspace{0.1cm}
		\textbf{OPTIONS:}
		\begin{enumerate}
			\item $\cos(\theta) = \frac{P + Q}{||P|| \times ||Q||}$
			\item $\cos(\theta) = \frac{P \cdot Q}{||P|| \times ||Q||}$
			\item $\cos(\theta) = \frac{||P||}{||Q||}$
			\item $\cos(\theta) = P \times Q$
		\end{enumerate}
		
		\vspace{0.1cm}
		\textcolor{myblue}{\textbf{ANSWER: B}}
		
		\vspace{0.1cm}
		\textbf{DETAILED EXPLANATION:}
		\begin{itemize}
			\item $P \cdot Q$ = dot product
			\item $||P||$ = magnitude (norm) of vector P
			\item Range: [-1, 1]
			\item 1 = identical, 0 = orthogonal, -1 = opposite
		\end{itemize}
		
		\vspace{0.1cm}
		\textcolor{myred}{\textbf{WHY OTHER OPTIONS ARE WRONG:}}
		\begin{itemize}
			\item \textcolor{myred}{A:} Addition is not the correct operation
			\item \textcolor{myred}{C:} Only compares magnitudes, ignores direction
			\item \textcolor{myred}{D:} Missing normalization (dividing by norms)
		\end{itemize}
		
		\textcolor{myred}{\textbf{EXAM TRAP:}} Cosine similarity = \textcolor{myblue}{dot product / (norm1 $\times$ norm2)}!
	\end{frame}
	
	\begin{frame}{Q19: Cosine Similarity Values - Tutorial}
		\fontsize{6.5}{7.5}\selectfont
		\textbf{QUESTION: What does a cosine similarity of 0.95 indicate?}
		
		\vspace{0.1cm}
		\textbf{OPTIONS:}
		\begin{enumerate}
			\item The words are completely unrelated
			\item The words are very similar (close in embedding space)
			\item The words are opposites
			\item The words are identical
		\end{enumerate}
		
		\vspace{0.1cm}
		\textcolor{myblue}{\textbf{ANSWER: B}}
		
		\vspace{0.1cm}
		\textbf{DETAILED EXPLANATION:}
		\begin{itemize}
			\item Range: -1 to 1
			\item 0.95 is very close to 1.0
			\item Indicates high similarity
			\item Words are likely similar in meaning
		\end{itemize}
		
		\vspace{0.1cm}
		\textcolor{myred}{\textbf{WHY OTHER OPTIONS ARE WRONG:}}
		\begin{itemize}
			\item \textcolor{myred}{A:} 0.95 is HIGH, not unrelated (0)
			\item \textcolor{myred}{C:} Opposites would be close to -1
			\item \textcolor{myred}{D:} Identical would be exactly 1.0
		\end{itemize}
		
		\textcolor{myred}{\textbf{EXAM TRAP:}} Higher cosine similarity = \textcolor{myblue}{more similar words}!
	\end{frame}
	
	\begin{frame}{Q20: Word Analogy Example - Tutorial}
		\fontsize{6.5}{7.5}\selectfont
		\textbf{QUESTION: What is the result of the operation "King - Man + Woman" in word embeddings?}
		
		\vspace{0.1cm}
		\textbf{OPTIONS:}
		\begin{enumerate}
			\item Prince
			\item Queen
			\item Royalty
			\item Knight
		\end{enumerate}
		
		\vspace{0.1cm}
		\textcolor{myblue}{\textbf{ANSWER: B}}
		
		\vspace{0.1cm}
		\textbf{DETAILED EXPLANATION:}
		\begin{itemize}
			\item "King - Man + Woman $\approx$ Queen"
			\item Captures gender relationship
			\item Vector arithmetic preserves semantic relationships
			\item Famous Word2Vec analogy
		\end{itemize}
		
		\vspace{0.1cm}
		\textcolor{myred}{\textbf{WHY OTHER OPTIONS ARE WRONG:}}
		\begin{itemize}
			\item \textcolor{myred}{A:} Prince is a different royal title
			\item \textcolor{myred}{C:} Royalty is too general
			\item \textcolor{myred}{D:} Knight is not a royal title
		\end{itemize}
		
		\textcolor{myred}{\textbf{EXAM TRAP:}} "King - Man + Woman" = \textcolor{myblue}{Queen}!
	\end{frame}
	
	\begin{frame}{Q21: Word Embedding Limitation - Tutorial}
		\fontsize{6.5}{7.5}\selectfont
		\textbf{QUESTION: What is a key limitation of word embeddings like Word2Vec?}
		
		\vspace{0.1cm}
		\textbf{OPTIONS:}
		\begin{enumerate}
			\item They capture sentence-level context perfectly
			\item They cannot capture context-dependent word meanings
			\item They always require a fixed context window
			\item Both B and C are limitations
		\end{enumerate}
		
		\vspace{0.1cm}
		\textcolor{myblue}{\textbf{ANSWER: D}}
		
		\vspace{0.1cm}
		\textbf{DETAILED EXPLANATION:}
		\begin{itemize}
			\item Same word has same embedding regardless of context
			\item "bat" (animal) vs "bat" (sports) - same embedding
			\item Fixed context window must be specified
			\item Cannot handle distant relationships beyond window
		\end{itemize}
		
		\vspace{0.1cm}
		\textcolor{myred}{\textbf{WHY OTHER OPTIONS ARE WRONG:}}
		\begin{itemize}
			\item \textcolor{myred}{A:} This is false - they DO NOT capture sentence-level context
			\item \textcolor{myred}{B:} This is true but incomplete
			\item \textcolor{myred}{C:} This is true but incomplete
		\end{itemize}
		
		\textcolor{myred}{\textbf{EXAM TRAP:}} Word2Vec gives \textcolor{myblue}{same vector for different contexts}!
	\end{frame}
	
	\begin{frame}{Q22: Other Embedding Approaches - Tutorial}
		\fontsize{6.5}{7.5}\selectfont
		\textbf{QUESTION: Which of the following is an alternative to Word2Vec for word embeddings?}
		
		\vspace{0.1cm}
		\textbf{OPTIONS:}
		\begin{enumerate}
			\item GloVe (Global Vectors)
			\item BERT (Bidirectional Encoder Representations)
			\item Both GloVe and BERT
			\item None of the above
		\end{enumerate}
		
		\vspace{0.1cm}
		\textcolor{myblue}{\textbf{ANSWER: C}}
		
		\vspace{0.1cm}
		\textbf{DETAILED EXPLANATION:}
		\begin{itemize}
			\item \textbf{GloVe:} Stanford University, combines matrix factorization
			\item \textbf{BERT:} Contextual embeddings (different vectors for different contexts)
			\item Both are alternatives/improvements over Word2Vec
		\end{itemize}
		
		\vspace{0.1cm}
		\textcolor{myred}{\textbf{WHY OTHER OPTIONS ARE WRONG:}}
		\begin{itemize}
			\item \textcolor{myred}{A:} GloVe is correct but incomplete
			\item \textcolor{myred}{B:} BERT is correct but incomplete
			\item \textcolor{myred}{D:} Both A and B are valid alternatives
		\end{itemize}
		
		\textcolor{myred}{\textbf{EXAM TRAP:}} GloVe and BERT are \textcolor{myblue}{both alternatives} to Word2Vec!
	\end{frame}
	
	\begin{frame}{Q23: t-SNE Purpose - Tutorial}
		\fontsize{6.5}{7.5}\selectfont
		\textbf{QUESTION: What is t-SNE used for in the context of word embeddings?}
		
		\vspace{0.1cm}
		\textbf{OPTIONS:}
		\begin{enumerate}
			\item To create word embeddings
			\item To visualize high-dimensional embeddings in 2D
			\item To calculate cosine similarity
			\item To train neural networks
		\end{enumerate}
		
		\vspace{0.1cm}
		\textcolor{myblue}{\textbf{ANSWER: B}}
		
		\vspace{0.1cm}
		\textbf{DETAILED EXPLANATION:}
		\begin{itemize}
			\item t-SNE = t-distributed Stochastic Neighbor Embedding
			\item Reduces high-dimensional embeddings to 2D
			\item Visualizes relationships between words
			\item Clusters show semantic groups
		\end{itemize}
		
		\vspace{0.1cm}
		\textcolor{myred}{\textbf{WHY OTHER OPTIONS ARE WRONG:}}
		\begin{itemize}
			\item \textcolor{myred}{A:} t-SNE visualizes, doesn't create embeddings
			\item \textcolor{myred}{C:} Cosine similarity is a separate measure
			\item \textcolor{myred}{D:} t-SNE is for visualization, not training
		\end{itemize}
		
		\textcolor{myred}{\textbf{EXAM TRAP:}} t-SNE is for \textcolor{myblue}{visualization}, not creation!
	\end{frame}
	
	\begin{frame}{Q24: Embedding Applications - Tutorial}
		\fontsize{6.5}{7.5}\selectfont
		\textbf{QUESTION: What is a common application of word embeddings?}
		
		\vspace{0.1cm}
		\textbf{OPTIONS:}
		\begin{enumerate}
			\item Sentiment analysis
			\item Document clustering
			\item Input to transformer models
			\item All of the above
		\end{enumerate}
		
		\vspace{0.1cm}
		\textcolor{myblue}{\textbf{ANSWER: D}}
		
		\vspace{0.1cm}
		\textbf{DETAILED EXPLANATION:}
		\begin{itemize}
			\item Clustering: group similar documents
			\item Sentiment analysis: understand word meanings
			\item Transformer input: pre-trained embeddings
			\item Foundation for modern NLP systems
		\end{itemize}
		
		\vspace{0.1cm}
		\textcolor{myred}{\textbf{WHY OTHER OPTIONS ARE WRONG:}}
		\begin{itemize}
			\item \textcolor{myred}{A:} Correct but incomplete
			\item \textcolor{myred}{B:} Correct but incomplete
			\item \textcolor{myred}{C:} Correct but incomplete
		\end{itemize}
		
		\textcolor{myred}{\textbf{EXAM TRAP:}} Embeddings are \textcolor{myblue}{foundational} for many NLP tasks!
	\end{frame}
	
	\begin{frame}{Q25: Sentence Embeddings - Tutorial}
		\fontsize{6.5}{7.5}\selectfont
		\textbf{QUESTION: What is SBert used for in NLP?}
		
		\vspace{0.1cm}
		\textbf{OPTIONS:}
		\begin{enumerate}
			\item Word-level embeddings
			\item Sentence-level embeddings
			\item Document-level embeddings
			\item Image generation
		\end{enumerate}
		
		\vspace{0.1cm}
		\textcolor{myblue}{\textbf{ANSWER: B}}
		
		\vspace{0.1cm}
		\textbf{DETAILED EXPLANATION:}
		\begin{itemize}
			\item SBert = Sentence BERT
			\item Captures sentence-level meaning
			\item Better for sentence comparison
			\item Beyond word-level embeddings
		\end{itemize}
		
		\vspace{0.1cm}
		\textcolor{myred}{\textbf{WHY OTHER OPTIONS ARE WRONG:}}
		\begin{itemize}
			\item \textcolor{myred}{A:} Word embeddings are Word2Vec/GloVe
			\item \textcolor{myred}{C:} Doc2Vec is for document embeddings
			\item \textcolor{myred}{D:} SBert is text-based, not image generation
		\end{itemize}
		
		\textcolor{myred}{\textbf{EXAM TRAP:}} SBert = \textcolor{myblue}{sentence-level embeddings}!
	\end{frame}
	
	% ============================================================
	% SECTION 3: RNNs AND LSTMs - QUESTIONS 26-35
	% ============================================================
	
	\begin{frame}{Q26: RNN Definition - Tutorial}
		\fontsize{6.5}{7.5}\selectfont
		\textbf{QUESTION: What is a Recurrent Neural Network (RNN)?}
		
		\vspace{0.1cm}
		\textbf{OPTIONS:}
		\begin{enumerate}
			\item A feedforward neural network with no loops
			\item A neural network designed for sequential data with memory
			\item A network that only processes images
			\item A network that has no hidden layers
		\end{enumerate}
		
		\vspace{0.1cm}
		\textcolor{myblue}{\textbf{ANSWER: B}}
		
		\vspace{0.1cm}
		\textbf{DETAILED EXPLANATION:}
		\begin{itemize}
			\item Has memory (hidden cell)
			\item Value at time t depends on previous values
			\item Designed for sequential data (text, speech, time series)
			\item Outputs become inputs for next step
		\end{itemize}
		
		\vspace{0.1cm}
		\textcolor{myred}{\textbf{WHY OTHER OPTIONS ARE WRONG:}}
		\begin{itemize}
			\item \textcolor{myred}{A:} FNNs have no loops; RNNs have feedback loops
			\item \textcolor{myred}{C:} RNNs handle sequential data, not just images
			\item \textcolor{myred}{D:} RNNs have hidden layers (that's their memory)
		\end{itemize}
		
		\textcolor{myred}{\textbf{EXAM TRAP:}} RNNs have \textcolor{myblue}{memory} for sequential data!
	\end{frame}
	
	\begin{frame}{Q27: RNN vs FNN - Tutorial}
		\fontsize{6.5}{7.5}\selectfont
		\textbf{QUESTION: What is the key difference between RNNs and Feedforward Neural Networks (FNNs)?}
		
		\vspace{0.1cm}
		\textbf{OPTIONS:}
		\begin{enumerate}
			\item FNNs have memory; RNNs do not
			\item RNNs have memory (hidden state); FNNs do not
			\item Both have memory
			\item Both are identical
		\end{enumerate}
		
		\vspace{0.1cm}
		\textcolor{myblue}{\textbf{ANSWER: B}}
		
		\vspace{0.1cm}
		\textbf{DETAILED EXPLANATION:}
		\begin{itemize}
			\item \textbf{FNN:} $y = f(x)$ - no memory
			\item \textbf{RNN:} $y_t = f(x_t, h_{t-1})$ - has memory
			\item RNNs process sequential data
			\item FNNs process static data
		\end{itemize}
		
		\vspace{0.1cm}
		\textcolor{myred}{\textbf{WHY OTHER OPTIONS ARE WRONG:}}
		\begin{itemize}
			\item \textcolor{myred}{A:} This reverses the definitions
			\item \textcolor{myred}{C:} FNNs do NOT have memory
			\item \textcolor{myred}{D:} They are fundamentally different
		\end{itemize}
		
		\textcolor{myred}{\textbf{EXAM TRAP:}} RNNs have \textcolor{myblue}{memory}; FNNs do not!
	\end{frame}
	
	\begin{frame}{Q28: RNN Architecture - Tutorial}
		\fontsize{6.5}{7.5}\selectfont
		\textbf{QUESTION: What is the hidden state $h_t$ in an RNN?}
		
		\vspace{0.1cm}
		\textbf{OPTIONS:}
		\begin{enumerate}
			\item The output at time t only
			\item The memory/state at time t that depends on previous states
			\item The input at time t
			\item The final output of the network
		\end{enumerate}
		
		\vspace{0.1cm}
		\textcolor{myblue}{\textbf{ANSWER: B}}
		
		\vspace{0.1cm}
		\textbf{DETAILED EXPLANATION:}
		\begin{itemize}
			\item $h_t = f(W_{xh}x_t + W_{hh}h_{t-1} + b_h)$
			\item Depends on current input AND previous state
			\item Stores information from previous time steps
			\item Represents the network's memory
		\end{itemize}
		
		\vspace{0.1cm}
		\textcolor{myred}{\textbf{WHY OTHER OPTIONS ARE WRONG:}}
		\begin{itemize}
			\item \textcolor{myred}{A:} Hidden state is not just the output
			\item \textcolor{myred}{C:} Input is $x_t$, not hidden state
			\item \textcolor{myred}{D:} Hidden state is intermediate, not final
		\end{itemize}
		
		\textcolor{myred}{\textbf{EXAM TRAP:}} Hidden state = \textcolor{myblue}{memory at time t}!
	\end{frame}
	
	\begin{frame}{Q29: Vanishing Gradient Problem - Tutorial}
		\fontsize{6.5}{7.5}\selectfont
		\textbf{QUESTION: What is the vanishing gradient problem in RNNs?}
		
		\vspace{0.1cm}
		\textbf{OPTIONS:}
		\begin{enumerate}
			\item Gradients become very large during training
			\item Gradients become very small during backpropagation, limiting long-range learning
			\item Gradients oscillate randomly
			\item Gradients remain constant throughout training
		\end{enumerate}
		
		\vspace{0.1cm}
		\textcolor{myblue}{\textbf{ANSWER: B}}
		
		\vspace{0.1cm}
		\textbf{DETAILED EXPLANATION:}
		\begin{itemize}
			\item Gradients multiplied through many time steps
			\item Each multiplication by $< 1$
			\item Product approaches zero
			\item Long-range dependencies are lost
		\end{itemize}
		
		\vspace{0.1cm}
		\textcolor{myred}{\textbf{WHY OTHER OPTIONS ARE WRONG:}}
		\begin{itemize}
			\item \textcolor{myred}{A:} This describes exploding gradients
			\item \textcolor{myred}{C:} This describes oscillations
			\item \textcolor{myred}{D:} Gradients change during training
		\end{itemize}
		
		\textcolor{myred}{\textbf{EXAM TRAP:}} Vanishing gradients = \textcolor{myblue}{very small gradients} in early layers!
	\end{frame}
	
	\begin{frame}{Q30: LSTM Purpose - Tutorial}
		\fontsize{6.5}{7.5}\selectfont
		\textbf{QUESTION: What was LSTM designed to address?}
		
		\vspace{0.1cm}
		\textbf{OPTIONS:}
		\begin{enumerate}
			\item The need for faster training
			\item The vanishing gradient problem in RNNs
			\item The need for larger datasets
			\item The need for simpler architectures
		\end{enumerate}
		
		\vspace{0.1cm}
		\textcolor{myblue}{\textbf{ANSWER: B}}
		
		\vspace{0.1cm}
		\textbf{DETAILED EXPLANATION:}
		\begin{itemize}
			\item LSTM = Long Short-Term Memory
			\item Specifically designed to address vanishing gradients
			\item Uses gates to regulate information flow
			\item Can capture long-range dependencies
		\end{itemize}
		
		\vspace{0.1cm}
		\textcolor{myred}{\textbf{WHY OTHER OPTIONS ARE WRONG:}}
		\begin{itemize}
			\item \textcolor{myred}{A:} LSTMs are more complex, not faster
			\item \textcolor{myred}{C:} LSTMs don't address data size issues
			\item \textcolor{myred}{D:} LSTMs are more complex, not simpler
		\end{itemize}
		
		\textcolor{myred}{\textbf{EXAM TRAP:}} LSTM = \textcolor{myblue}{solution to vanishing gradient} problem!
	\end{frame}
	
	\begin{frame}{Q31: LSTM Gates - Tutorial}
		\fontsize{6.5}{7.5}\selectfont
		\textbf{QUESTION: What are the three gates in a basic LSTM cell?}
		
		\vspace{0.1cm}
		\textbf{OPTIONS:}
		\begin{enumerate}
			\item Input, Output, and Hidden
			\item Forget, Input, and Output
			\item Memory, Forget, and Output
			\item Encoder, Decoder, and Memory
		\end{enumerate}
		
		\vspace{0.1cm}
		\textcolor{myblue}{\textbf{ANSWER: B}}
		
		\vspace{0.1cm}
		\textbf{DETAILED EXPLANATION:}
		\begin{itemize}
			\item \textbf{Forget Gate:} What to discard from memory
			\item \textbf{Input Gate:} What new information to store
			\item \textbf{Output Gate:} What to output
			\item These gates regulate information flow
		\end{itemize}
		
		\vspace{0.1cm}
		\textcolor{myred}{\textbf{WHY OTHER OPTIONS ARE WRONG:}}
		\begin{itemize}
			\item \textcolor{myred}{A:} "Hidden" is not a gate
			\item \textcolor{myred}{C:} "Memory" is not a gate
			\item \textcolor{myred}{D:} Encoder/Decoder are for transformers, not LSTMs
		\end{itemize}
		
		\textcolor{myred}{\textbf{EXAM TRAP:}} LSTM gates = \textcolor{myblue}{Forget, Input, Output}!
	\end{frame}
	
	\begin{frame}{Q32: LSTM vs RNN Comparison - Tutorial}
		\fontsize{6.5}{7.5}\selectfont
		\textbf{QUESTION: Which statement correctly compares LSTM and RNN?}
		
		\vspace{0.1cm}
		\textbf{OPTIONS:}
		\begin{enumerate}
			\item RNN is more powerful than LSTM
			\item LSTM handles long-term dependencies better than RNN
			\item RNN and LSTM are identical
			\item LSTM is simpler than RNN
		\end{enumerate}
		
		\vspace{0.1cm}
		\textcolor{myblue}{\textbf{ANSWER: B}}
		
		\vspace{0.1cm}
		\textbf{DETAILED EXPLANATION:}
		\begin{itemize}
			\item RNN struggles with long-term dependencies (vanishing gradients)
			\item LSTM solves this with gates
			\item LSTM is more powerful but more complex
			\item LSTM can capture relationships over longer distances
		\end{itemize}
		
		\vspace{0.1cm}
		\textcolor{myred}{\textbf{WHY OTHER OPTIONS ARE WRONG:}}
		\begin{itemize}
			\item \textcolor{myred}{A:} LSTM is more powerful than RNN
			\item \textcolor{myred}{C:} They are different architectures
			\item \textcolor{myred}{D:} LSTM is more complex, not simpler
		\end{itemize}
		
		\textcolor{myred}{\textbf{EXAM TRAP:}} LSTM \textcolor{myblue}{outperforms} RNN for long sequences!
	\end{frame}
	
	\begin{frame}{Q33: LSTM Forget Gate - Tutorial}
		\fontsize{6.5}{7.5}\selectfont
		\textbf{QUESTION: What is the purpose of the Forget Gate in an LSTM?}
		
		\vspace{0.1cm}
		\textbf{OPTIONS:}
		\begin{enumerate}
			\item To determine what new information to store
			\item To determine what information to discard from memory
			\item To determine what to output
			\item To initialize the network
		\end{enumerate}
		
		\vspace{0.1cm}
		\textcolor{myblue}{\textbf{ANSWER: B}}
		
		\vspace{0.1cm}
		\textbf{DETAILED EXPLANATION:}
		\begin{itemize}
			\item $f_t = \sigma(W_f \cdot [h_{t-1}, x_t] + b_f)$
			\item Decides which information is no longer relevant
			\item Prevents memory from overflowing
			\item Essential for long-term learning
		\end{itemize}
		
		\vspace{0.1cm}
		\textcolor{myred}{\textbf{WHY OTHER OPTIONS ARE WRONG:}}
		\begin{itemize}
			\item \textcolor{myred}{A:} This is the Input Gate's function
			\item \textcolor{myred}{C:} This is the Output Gate's function
			\item \textcolor{myred}{D:} This is initialization, not a gate function
		\end{itemize}
		
		\textcolor{myred}{\textbf{EXAM TRAP:}} Forget Gate = \textcolor{myblue}{what to discard from memory}!
	\end{frame}
	
	\begin{frame}{Q34: RNN Applications - Tutorial}
		\fontsize{6.5}{7.5}\selectfont
		\textbf{QUESTION: What is a common application of RNNs in finance?}
		
		\vspace{0.1cm}
		\textbf{OPTIONS:}
		\begin{enumerate}
			\item Image classification
			\item Time series analysis and prediction
			\item Image generation
			\item Image segmentation
		\end{enumerate}
		
		\vspace{0.1cm}
		\textcolor{myblue}{\textbf{ANSWER: B}}
		
		\vspace{0.1cm}
		\textbf{DETAILED EXPLANATION:}
		\begin{itemize}
			\item RNNs excel at sequential data
			\item Stock price prediction
			\item Volatility forecasting
			\item Economic indicator forecasting
		\end{itemize}
		
		\vspace{0.1cm}
		\textcolor{myred}{\textbf{WHY OTHER OPTIONS ARE WRONG:}}
		\begin{itemize}
			\item \textcolor{myred}{A:} CNNs are better for image classification
			\item \textcolor{myred}{C:} GANs/Diffusion models are for image generation
			\item \textcolor{myred}{D:} CNNs are used for image segmentation
		\end{itemize}
		
		\textcolor{myred}{\textbf{EXAM TRAP:}} RNNs = \textcolor{myblue}{time series and sequential data}!
	\end{frame}
	
	\begin{frame}{Q35: RNN Limitations - Tutorial}
		\fontsize{6.5}{7.5}\selectfont
		\textbf{QUESTION: What is a major limitation of RNNs that transformers address?}
		
		\vspace{0.1cm}
		\textbf{OPTIONS:}
		\begin{enumerate}
			\item RNNs are too fast
			\item RNNs cannot parallelize effectively
			\item RNNs require too little data
			\item RNNs are too interpretable
		\end{enumerate}
		
		\vspace{0.1cm}
		\textcolor{myblue}{\textbf{ANSWER: B}}
		
		\vspace{0.1cm}
		\textbf{DETAILED EXPLANATION:}
		\begin{itemize}
			\item RNNs process sequentially (word by word)
			\item Cannot parallelize across sequence
			\item Transformers process entire sequence in parallel
			\item Transformers are much faster
		\end{itemize}
		
		\vspace{0.1cm}
		\textcolor{myred}{\textbf{WHY OTHER OPTIONS ARE WRONG:}}
		\begin{itemize}
			\item \textcolor{myred}{A:} RNNs are slow, not too fast
			\item \textcolor{myred}{C:} RNNs require lots of data
			\item \textcolor{myred}{D:} RNNs are harder to interpret than feedforward, not more
		\end{itemize}
		
		\textcolor{myred}{\textbf{EXAM TRAP:}} RNNs cannot \textcolor{myblue}{parallelize}; transformers can!
	\end{frame}
	
	% ============================================================
	% SECTION 4: TRANSFORMERS - QUESTIONS 36-50
	% ============================================================
	
	\begin{frame}{Q36: Transformer Paper - Tutorial}
		\fontsize{6.5}{7.5}\selectfont
		\textbf{QUESTION: What paper introduced the transformer architecture?}
		
		\vspace{0.1cm}
		\textbf{OPTIONS:}
		\begin{enumerate}
			\item "Attention is All You Need" (2017)
			\item "BERT: Pre-training of Deep Bidirectional Transformers" (2018)
			\item "Generative Pre-trained Transformer" (2018)
			\item "Language Models are Few-Shot Learners" (2020)
		\end{enumerate}
		
		\vspace{0.1cm}
		\textcolor{myblue}{\textbf{ANSWER: A}}
		
		\vspace{0.1cm}
		\textbf{DETAILED EXPLANATION:}
		\begin{itemize}
			\item "Attention is All You Need" by Vaswani et al. (2017)
			\item Introduced the transformer architecture
			\item Revolutionized NLP
			\item Foundation for all modern LLMs
		\end{itemize}
		
		\vspace{0.1cm}
		\textcolor{myred}{\textbf{WHY OTHER OPTIONS ARE WRONG:}}
		\begin{itemize}
			\item \textcolor{myred}{B:} BERT came after transformers
			\item \textcolor{myred}{C:} GPT came after transformers
			\item \textcolor{myred}{D:} GPT-3 paper came after transformers
		\end{itemize}
		
		\textcolor{myred}{\textbf{EXAM TRAP:}} Transformers introduced in \textcolor{myblue}{"Attention is All You Need"} (2017)!
	\end{frame}
	
	\begin{frame}{Q37: Transformer Components - Tutorial}
		\fontsize{6.5}{7.5}\selectfont
		\textbf{QUESTION: What are the two main components of a transformer?}
		
		\vspace{0.1cm}
		\textbf{OPTIONS:}
		\begin{enumerate}
			\item Generator and Discriminator
			\item Encoder and Decoder
			\item Input and Output layers only
			\item Attention and Recurrent layers
		\end{enumerate}
		
		\vspace{0.1cm}
		\textcolor{myblue}{\textbf{ANSWER: B}}
		
		\vspace{0.1cm}
		\textbf{DETAILED EXPLANATION:}
		\begin{itemize}
			\item \textbf{Encoder:} Processes input, converts to vectors, "understands" text
			\item \textbf{Decoder:} Generates output based on input/prompt
			\item Some models use only encoder (BERT) or only decoder (GPT)
			\item Original had both (encoder-decoder)
		\end{itemize}
		
		\vspace{0.1cm}
		\textcolor{myred}{\textbf{WHY OTHER OPTIONS ARE WRONG:}}
		\begin{itemize}
			\item \textcolor{myred}{A:} These are GAN components, not transformer
			\item \textcolor{myred}{C:} Transformers have more than just input/output
			\item \textcolor{myred}{D:} Transformers use attention, not recurrence
		\end{itemize}
		
		\textcolor{myred}{\textbf{EXAM TRAP:}} Transformer = \textcolor{myblue}{Encoder + Decoder}!
	\end{frame}
	
	\begin{frame}{Q38: Self-Attention Definition - Tutorial}
		\fontsize{6.5}{7.5}\selectfont
		\textbf{QUESTION: What is self-attention in a transformer?}
		
		\vspace{0.1cm}
		\textbf{OPTIONS:}
		\begin{enumerate}
			\item Ignoring all other words in a sentence
			\item Calculating relationships between all words in a sequence
			\item Only paying attention to the first word
			\item Paying attention only to nouns
		\end{enumerate}
		
		\vspace{0.1cm}
		\textcolor{myblue}{\textbf{ANSWER: B}}
		
		\vspace{0.1cm}
		\textbf{DETAILED EXPLANATION:}
		\begin{itemize}
			\item Each word analyzes all other words
			\item Captures relationships between all words
			\item Creates semantic understanding
			\item Enables long-range dependencies
		\end{itemize}
		
		\vspace{0.1cm}
		\textcolor{myred}{\textbf{WHY OTHER OPTIONS ARE WRONG:}}
		\begin{itemize}
			\item \textcolor{myred}{A:} This would be no-attention
			\item \textcolor{myred}{C:} This ignores important context
			\item \textcolor{myred}{D:} Attention applies to ALL words
		\end{itemize}
		
		\textcolor{myred}{\textbf{EXAM TRAP:}} Self-attention = \textcolor{myblue}{all words analyze all other words}!
	\end{frame}
	
	\begin{frame}{Q39: Query, Key, Value - Tutorial}
		\fontsize{6.5}{7.5}\selectfont
		\textbf{QUESTION: What are the three roles in the attention mechanism?}
		
		\vspace{0.1cm}
		\textbf{OPTIONS:}
		\begin{enumerate}
			\item Input, Output, and Hidden
			\item Query, Key, and Value
			\item Encoder, Decoder, and Memory
			\item Training, Validation, and Testing
		\end{enumerate}
		
		\vspace{0.1cm}
		\textcolor{myblue}{\textbf{ANSWER: B}}
		
		\vspace{0.1cm}
		\textbf{DETAILED EXPLANATION:}
		\begin{itemize}
			\item \textbf{Query (Q):} Current token asking "what should I attend to?"
			\item \textbf{Key (K):} Other tokens responding to the query
			\item \textbf{Value (V):} Content of tokens used to create context
			\item All three are derived from input embeddings
		\end{itemize}
		
		\vspace{0.1cm}
		\textcolor{myred}{\textbf{WHY OTHER OPTIONS ARE WRONG:}}
		\begin{itemize}
			\item \textcolor{myred}{A:} These are general terms, not attention-specific
			\item \textcolor{myred}{C:} Encoder/Decoder are transformer components
			\item \textcolor{myred}{D:} These are ML phases, not attention roles
		\end{itemize}
		
		\textcolor{myred}{\textbf{EXAM TRAP:}} Attention uses \textcolor{myblue}{Query, Key, Value} vectors!
	\end{frame}
	
	\begin{frame}{Q40: Attention Score Calculation - Tutorial}
		\fontsize{6.5}{7.5}\selectfont
		\textbf{QUESTION: How is the attention score between query $q_i$ and key $k_j$ calculated?}
		
		\vspace{0.1cm}
		\textbf{OPTIONS:}
		\begin{enumerate}
			\item $score(q_i, k_j) = q_i + k_j$
			\item $score(q_i, k_j) = q_i \cdot k_j$
			\item $score(q_i, k_j) = q_i - k_j$
			\item $score(q_i, k_j) = q_i / k_j$
		\end{enumerate}
		
		\vspace{0.1cm}
		\textcolor{myblue}{\textbf{ANSWER: B}}
		
		\vspace{0.1cm}
		\textbf{DETAILED EXPLANATION:}
		\begin{itemize}
			\item Dot product measures similarity
			\item Higher dot product = more relevant
			\item Then normalized with softmax
			\item Weighted sum of value vectors
		\end{itemize}
		
		\vspace{0.1cm}
		\textcolor{myred}{\textbf{WHY OTHER OPTIONS ARE WRONG:}}
		\begin{itemize}
			\item \textcolor{myred}{A:} Addition doesn't measure similarity
			\item \textcolor{myred}{C:} Subtraction doesn't measure similarity
			\item \textcolor{myred}{D:} Division doesn't measure similarity
		\end{itemize}
		
		\textcolor{myred}{\textbf{EXAM TRAP:}} Attention score = \textcolor{myblue}{dot product of Query and Key}!
	\end{frame}
	
	\begin{frame}{Q41: Multi-Head Attention - Tutorial}
		\fontsize{6.5}{7.5}\selectfont
		\textbf{QUESTION: What is the purpose of multi-head attention?}
		
		\vspace{0.1cm}
		\textbf{OPTIONS:}
		\begin{enumerate}
			\item To reduce computation
			\item To capture different types of relationships in parallel
			\item To make the model simpler
			\item To eliminate the need for feedforward layers
		\end{enumerate}
		
		\vspace{0.1cm}
		\textcolor{myblue}{\textbf{ANSWER: B}}
		
		\vspace{0.1cm}
		\textbf{DETAILED EXPLANATION:}
		\begin{itemize}
			\item Multiple attention heads in parallel
			\item Each head focuses on different aspects
			\item Captures different relationships
			\item Example: GPT has 12 heads
		\end{itemize}
		
		\vspace{0.1cm}
		\textcolor{myred}{\textbf{WHY OTHER OPTIONS ARE WRONG:}}
		\begin{itemize}
			\item \textcolor{myred}{A:} Multi-head increases computation
			\item \textcolor{myred}{C:} Multi-head makes models more complex
			\item \textcolor{myred}{D:} Feedforward layers are still needed
		\end{itemize}
		
		\textcolor{myred}{\textbf{EXAM TRAP:}} Multi-head captures \textcolor{myblue}{diverse relationships}!
	\end{frame}
	
	\begin{frame}{Q42: Positional Embeddings - Tutorial}
		\fontsize{6.5}{7.5}\selectfont
		\textbf{QUESTION: Why are positional embeddings needed in transformers?}
		
		\vspace{0.1cm}
		\textbf{OPTIONS:}
		\begin{enumerate}
			\item Transformers process words in order, so they don't need positions
			\item Transformers process words in parallel, so they need position information
			\item Positional embeddings are optional
			\item Positional embeddings are only for RNNs
		\end{enumerate}
		
		\vspace{0.1cm}
		\textcolor{myblue}{\textbf{ANSWER: B}}
		
		\vspace{0.1cm}
		\textbf{DETAILED EXPLANATION:}
		\begin{itemize}
			\item Transformers process all words simultaneously
			\item Without positional info, word order is lost
			\item Positional embeddings encode position
			\item "dog bites man" vs "man bites dog"
		\end{itemize}
		
		\vspace{0.1cm}
		\textcolor{myred}{\textbf{WHY OTHER OPTIONS ARE WRONG:}}
		\begin{itemize}
			\item \textcolor{myred}{A:} This is incorrect - positions are needed
			\item \textcolor{myred}{C:} Positional embeddings are necessary
			\item \textcolor{myred}{D:} Positional embeddings are for transformers
		\end{itemize}
		
		\textcolor{myred}{\textbf{EXAM TRAP:}} Positional embeddings = \textcolor{myblue}{needed because transformers process in parallel}!
	\end{frame}
	
	\begin{frame}{Q43: Transformer Advantage - Tutorial}
		\fontsize{6.5}{7.5}\selectfont
		\textbf{QUESTION: What is a key advantage of transformers over RNNs?}
		
		\vspace{0.1cm}
		\textbf{OPTIONS:}
		\begin{enumerate}
			\item Transformers are always more interpretable
			\item Transformers process in parallel, making them faster
			\item Transformers require less data
			\item Transformers are simpler
		\end{enumerate}
		
		\vspace{0.1cm}
		\textcolor{myblue}{\textbf{ANSWER: B}}
		
		\vspace{0.1cm}
		\textbf{DETAILED EXPLANATION:}
		\begin{itemize}
			\item RNNs process sequentially (slow)
			\item Transformers process entire sequence at once (fast)
			\item Parallel processing = faster training
			\item Better hardware utilization
		\end{itemize}
		
		\vspace{0.1cm}
		\textcolor{myred}{\textbf{WHY OTHER OPTIONS ARE WRONG:}}
		\begin{itemize}
			\item \textcolor{myred}{A:} Transformers are LESS interpretable than RNNs
			\item \textcolor{myred}{C:} Transformers require MORE data
			\item \textcolor{myred}{D:} Transformers are MORE complex
		\end{itemize}
		
		\textcolor{myred}{\textbf{EXAM TRAP:}} Transformers = \textcolor{myblue}{parallel processing} = faster!
	\end{frame}
	
	\begin{frame}{Q44: BERT Architecture - Tutorial}
		\fontsize{6.5}{7.5}\selectfont
		\textbf{QUESTION: What type of architecture does BERT use?}
		
		\vspace{0.1cm}
		\textbf{OPTIONS:}
		\begin{enumerate}
			\item Encoder-only
			\item Decoder-only
			\item Encoder-Decoder
			\item Neither encoder nor decoder
		\end{enumerate}
		
		\vspace{0.1cm}
		\textcolor{myblue}{\textbf{ANSWER: A}}
		
		\vspace{0.1cm}
		\textbf{DETAILED EXPLANATION:}
		\begin{itemize}
			\item BERT = Bidirectional Encoder Representations
			\item Uses only the encoder part
			\item No decoder (not generative)
			\item Good for understanding tasks
		\end{itemize}
		
		\vspace{0.1cm}
		\textcolor{myred}{\textbf{WHY OTHER OPTIONS ARE WRONG:}}
		\begin{itemize}
			\item \textcolor{myred}{B:} GPT uses decoder-only
			\item \textcolor{myred}{C:} BART uses encoder-decoder
			\item \textcolor{myred}{D:} BERT uses encoders
		\end{itemize}
		
		\textcolor{myred}{\textbf{EXAM TRAP:}} BERT = \textcolor{myblue}{encoder-only} architecture!
	\end{frame}
	
	\begin{frame}{Q45: BERT Training - Tutorial}
		\fontsize{6.5}{7.5}\selectfont
		\textbf{QUESTION: How is BERT trained?}
		
		\vspace{0.1cm}
		\textbf{OPTIONS:}
		\begin{enumerate}
			\item Only using labeled data
			\item Using masked language modeling and next sentence prediction
			\item Using reinforcement learning
			\item Using only supervised learning
		\end{enumerate}
		
		\vspace{0.1cm}
		\textcolor{myblue}{\textbf{ANSWER: B}}
		
		\vspace{0.1cm}
		\textbf{DETAILED EXPLANATION:}
		\begin{itemize}
			\item \textbf{Masked Language Modeling:} 15\% of words masked, model predicts them
			\item \textbf{Next Sentence Prediction:} Predict if sentences are sequential
			\item Self-supervised (labels come from data)
			\item Learns bidirectional context
		\end{itemize}
		
		\vspace{0.1cm}
		\textcolor{myred}{\textbf{WHY OTHER OPTIONS ARE WRONG:}}
		\begin{itemize}
			\item \textcolor{myred}{A:} BERT uses unlabeled data for pre-training
			\item \textcolor{myred}{C:} Reinforcement learning is for RLHF
			\item \textcolor{myred}{D:} BERT uses self-supervised learning
		\end{itemize}
		
		\textcolor{myred}{\textbf{EXAM TRAP:}} BERT = \textcolor{myblue}{masked LM + next sentence prediction}!
	\end{frame}
	
	\begin{frame}{Q46: GPT Architecture - Tutorial}
		\fontsize{6.5}{7.5}\selectfont
		\textbf{QUESTION: What type of architecture does GPT use?}
		
		\vspace{0.1cm}
		\textbf{OPTIONS:}
		\begin{enumerate}
			\item Encoder-only
			\item Decoder-only
			\item Encoder-Decoder
			\item No transformer architecture
		\end{enumerate}
		
		\vspace{0.1cm}
		\textcolor{myblue}{\textbf{ANSWER: B}}
		
		\vspace{0.1cm}
		\textbf{DETAILED EXPLANATION:}
		\begin{itemize}
			\item GPT = Generative Pre-trained Transformer
			\item Uses only the decoder part
			\item Autoregressive (predicts next token)
			\item Good for generation tasks
		\end{itemize}
		
		\vspace{0.1cm}
		\textcolor{myred}{\textbf{WHY OTHER OPTIONS ARE WRONG:}}
		\begin{itemize}
			\item \textcolor{myred}{A:} BERT uses encoder-only
			\item \textcolor{myred}{C:} BART uses encoder-decoder
			\item \textcolor{myred}{D:} GPT uses transformer architecture
		\end{itemize}
		
		\textcolor{myred}{\textbf{EXAM TRAP:}} GPT = \textcolor{myblue}{decoder-only} architecture!
	\end{frame}
	
	\begin{frame}{Q47: BART Architecture - Tutorial}
		\fontsize{6.5}{7.5}\selectfont
		\textbf{QUESTION: What type of architecture does BART use?}
		
		\vspace{0.1cm}
		\textbf{OPTIONS:}
		\begin{enumerate}
			\item Encoder-only
			\item Decoder-only
			\item Encoder-Decoder
			\item No transformer architecture
		\end{enumerate}
		
		\vspace{0.1cm}
		\textcolor{myblue}{\textbf{ANSWER: C}}
		
		\vspace{0.1cm}
		\textbf{DETAILED EXPLANATION:}
		\begin{itemize}
			\item BART = Bidirectional and Auto-Regressive Transformer
			\item Uses both encoder and decoder
			\item Most versatile
			\item Good for many tasks (summarization, translation)
		\end{itemize}
		
		\vspace{0.1cm}
		\textcolor{myred}{\textbf{WHY OTHER OPTIONS ARE WRONG:}}
		\begin{itemize}
			\item \textcolor{myred}{A:} BERT uses encoder-only
			\item \textcolor{myred}{B:} GPT uses decoder-only
			\item \textcolor{myred}{D:} BART uses transformer architecture
		\end{itemize}
		
		\textcolor{myred}{\textbf{EXAM TRAP:}} BART = \textcolor{myblue}{encoder-decoder} architecture!
	\end{frame}
	
	\begin{frame}{Q48: RoBERTa vs BERT - Tutorial}
		\fontsize{6.5}{7.5}\selectfont
		\textbf{QUESTION: How does RoBERTa improve upon BERT?}
		
		\vspace{0.1cm}
		\textbf{OPTIONS:}
		\begin{enumerate}
			\item It uses fewer training cycles
			\item It uses different masking patterns and larger dataset
			\item It removes the encoder
			\item It uses only supervised learning
		\end{enumerate}
		
		\vspace{0.1cm}
		\textcolor{myblue}{\textbf{ANSWER: B}}
		
		\vspace{0.1cm}
		\textbf{DETAILED EXPLANATION:}
		\begin{itemize}
			\item RoBERTa = Robustly Optimized BERT Approach
			\item Different masking patterns each cycle
			\item Larger training dataset
			\item More training cycles
			\item Multi-language web crawled data
		\end{itemize}
		
		\vspace{0.1cm}
		\textcolor{myred}{\textbf{WHY OTHER OPTIONS ARE WRONG:}}
		\begin{itemize}
			\item \textcolor{myred}{A:} RoBERTa uses MORE training cycles
			\item \textcolor{myred}{C:} RoBERTa still uses encoders
			\item \textcolor{myred}{D:} RoBERTa uses self-supervised learning
		\end{itemize}
		
		\textcolor{myred}{\textbf{EXAM TRAP:}} RoBERTa = \textcolor{myblue}{improved BERT} with more data and better masking!
	\end{frame}
	
	\begin{frame}{Q49: FinBERT Purpose - Tutorial}
		\fontsize{6.5}{7.5}\selectfont
		\textbf{QUESTION: What is FinBERT specifically designed for?}
		
		\vspace{0.1cm}
		\textbf{OPTIONS:}
		\begin{enumerate}
			\item General purpose NLP
			\item Financial sentiment analysis
			\item Image generation
			\item Speech recognition
		\end{enumerate}
		
		\vspace{0.1cm}
		\textcolor{myblue}{\textbf{ANSWER: B}}
		
		\vspace{0.1cm}
		\textbf{DETAILED EXPLANATION:}
		\begin{itemize}
			\item FinBERT is a financial text specialization
			\item Trained on Reuters financial documents
			\item Corporate statements and analyst reports
			\item Optimal for financial sentiment analysis
		\end{itemize}
		
		\vspace{0.1cm}
		\textcolor{myred}{\textbf{WHY OTHER OPTIONS ARE WRONG:}}
		\begin{itemize}
			\item \textcolor{myred}{A:} FinBERT is domain-specific, not general
			\item \textcolor{myred}{C:} FinBERT is for text, not images
			\item \textcolor{myred}{D:} FinBERT is for text analysis, not speech
		\end{itemize}
		
		\textcolor{myred}{\textbf{EXAM TRAP:}} FinBERT = \textcolor{myblue}{financial sentiment analysis}!
	\end{frame}
	
	\begin{frame}{Q50: Transformer Disadvantages - Tutorial}
		\fontsize{6.5}{7.5}\selectfont
		\textbf{QUESTION: What is a disadvantage of transformers?}
		
		\vspace{0.1cm}
		\textbf{OPTIONS:}
		\begin{enumerate}
			\item They are too interpretable
			\item They are computationally intensive
			\item They cannot handle long-range dependencies
			\item They process sequentially
		\end{enumerate}
		
		\vspace{0.1cm}
		\textcolor{myblue}{\textbf{ANSWER: B}}
		
		\vspace{0.1cm}
		\textbf{DETAILED EXPLANATION:}
		\begin{itemize}
			\item Transformers require powerful hardware
			\item More computationally intensive than RNNs
			\item Higher energy consumption
			\item Longer training times
		\end{itemize}
		
		\vspace{0.1cm}
		\textcolor{myred}{\textbf{WHY OTHER OPTIONS ARE WRONG:}}
		\begin{itemize}
			\item \textcolor{myred}{A:} Transformers are LESS interpretable
			\item \textcolor{myred}{C:} Transformers EXCEL at long-range dependencies
			\item \textcolor{myred}{D:} Transformers process in parallel, not sequentially
		\end{itemize}
		
		\textcolor{myred}{\textbf{EXAM TRAP:}} Transformers are \textcolor{myblue}{resource-intensive}!
	\end{frame}
	
	% ============================================================
	% SECTION 5: LARGE LANGUAGE MODELS - QUESTIONS 51-70
	% ============================================================
	
	\begin{frame}{Q51: LLM Parameter Count - Tutorial}
		\fontsize{6.5}{7.5}\selectfont
		\textbf{QUESTION: Approximately how many parameters does GPT-3 have?}
		
		\vspace{0.1cm}
		\textbf{OPTIONS:}
		\begin{enumerate}
			\item 100 million
			\item 175 billion
			\item 1 trillion
			\item 1 million
		\end{enumerate}
		
		\vspace{0.1cm}
		\textcolor{myblue}{\textbf{ANSWER: B}}
		
		\vspace{0.1cm}
		\textbf{DETAILED EXPLANATION:}
		\begin{itemize}
			\item GPT-3 has approximately 175 billion parameters
			\item Much larger than GPT-1 (100 million)
			\item Enables generative capacity
			\item Trained on vast web text
		\end{itemize}
		
		\vspace{0.1cm}
		\textcolor{myred}{\textbf{WHY OTHER OPTIONS ARE WRONG:}}
		\begin{itemize}
			\item \textcolor{myred}{A:} This is GPT-1's size
			\item \textcolor{myred}{C:} Some newer models approach this, but not GPT-3
			\item \textcolor{myred}{D:} This is far too small
		\end{itemize}
		
		\textcolor{myred}{\textbf{EXAM TRAP:}} GPT-3 = \textcolor{myblue}{175 billion} parameters!
	\end{frame}
	
	\begin{frame}{Q52: LLM Types - Tutorial}
		\fontsize{6.5}{7.5}\selectfont
		\textbf{QUESTION: What are the three types of LLMs based on architecture?}
		
		\vspace{0.1cm}
		\textbf{OPTIONS:}
		\begin{enumerate}
			\item Autoencoding, Autoregressive, and Combination
			\item Encoder, Decoder, and Memory
			\item Training, Validation, and Testing
			\item Supervised, Unsupervised, and Reinforcement
		\end{enumerate}
		
		\vspace{0.1cm}
		\textcolor{myblue}{\textbf{ANSWER: A}}
		
		\vspace{0.1cm}
		\textbf{DETAILED EXPLANATION:}
		\begin{itemize}
			\item \textbf{Autoencoding:} BERT (predicts masked words)
			\item \textbf{Autoregressive:} GPT (predicts next word)
			\item \textbf{Combination:} BART (both encoder and decoder)
		\end{itemize}
		
		\vspace{0.1cm}
		\textcolor{myred}{\textbf{WHY OTHER OPTIONS ARE WRONG:}}
		\begin{itemize}
			\item \textcolor{myred}{B:} These are components, not LLM types
			\item \textcolor{myred}{C:} These are ML phases, not LLM types
			\item \textcolor{myred}{D:} These are learning paradigms, not LLM types
		\end{itemize}
		
		\textcolor{myred}{\textbf{EXAM TRAP:}} LLM types = \textcolor{myblue}{Autoencoding, Autoregressive, Combination}!
	\end{frame}
	
	\begin{frame}{Q53: LLM Training Step 1 - Tutorial}
		\fontsize{6.5}{7.5}\selectfont
		\textbf{QUESTION: What is the first step in LLM training?}
		
		\vspace{0.1cm}
		\textbf{OPTIONS:}
		\begin{enumerate}
			\item Fine-tuning
			\item Data preparation (cleaning, filtering, balancing)
			\item Decoding
			\item Reinforcement learning
		\end{enumerate}
		
		\vspace{0.1cm}
		\textcolor{myblue}{\textbf{ANSWER: B}}
		
		\vspace{0.1cm}
		\textbf{DETAILED EXPLANATION:}
		\begin{itemize}
			\item Clean training data (corpus)
			\item Filter and remove duplicates
			\item Remove noisy data and punctuation
			\item Balance class distributions
			\item Foundation for all subsequent steps
		\end{itemize}
		
		\vspace{0.1cm}
		\textcolor{myred}{\textbf{WHY OTHER OPTIONS ARE WRONG:}}
		\begin{itemize}
			\item \textcolor{myred}{A:} Fine-tuning happens after pre-training
			\item \textcolor{myred}{C:} Decoding happens at the end
			\item \textcolor{myred}{D:} RLHF happens during alignment
		\end{itemize}
		
		\textcolor{myred}{\textbf{EXAM TRAP:}} First step = \textcolor{myblue}{data preparation}!
	\end{frame}
	
	\begin{frame}{Q54: Tokenization in LLMs - Tutorial}
		\fontsize{6.5}{7.5}\selectfont
		\textbf{QUESTION: What is tokenization in LLM training?}
		
		\vspace{0.1cm}
		\textbf{OPTIONS:}
		\begin{enumerate}
			\item Removing all punctuation
			\item Separating data into small parts (tokens) for efficient processing
			\item Converting text to lowercase
			\item Removing stop words
		\end{enumerate}
		
		\vspace{0.1cm}
		\textcolor{myblue}{\textbf{ANSWER: B}}
		
		\vspace{0.1cm}
		\textbf{DETAILED EXPLANATION:}
		\begin{itemize}
			\item Split text into tokens (words or subwords)
			\item Enables efficient processing
			\item Store positional information
			\item Foundation for model understanding
		\end{itemize}
		
		\vspace{0.1cm}
		\textcolor{myred}{\textbf{WHY OTHER OPTIONS ARE WRONG:}}
		\begin{itemize}
			\item \textcolor{myred}{A:} Tokenization is more than removing punctuation
			\item \textcolor{myred}{C:} Lowercasing is a pre-processing step
			\item \textcolor{myred}{D:} Stop word removal is a pre-processing step
		\end{itemize}
		
		\textcolor{myred}{\textbf{EXAM TRAP:}} Tokenization = \textcolor{myblue}{splitting into tokens}!
	\end{frame}
	
	\begin{frame}{Q55: Pre-training in LLMs - Tutorial}
		\fontsize{6.5}{7.5}\selectfont
		\textbf{QUESTION: What happens during LLM pre-training?}
		
		\vspace{0.1cm}
		\textbf{OPTIONS:}
		\begin{enumerate}
			\item The model is fine-tuned on labeled data
			\item The model is trained on unlabeled text to understand word relationships
			\item The model is deployed in production
			\item The model is evaluated on test data
		\end{enumerate}
		
		\vspace{0.1cm}
		\textcolor{myblue}{\textbf{ANSWER: B}}
		
		\vspace{0.1cm}
		\textbf{DETAILED EXPLANATION:}
		\begin{itemize}
			\item Trained on unlabeled text
			\item Understands word relationships
			\item Learns language patterns
			\item Foundation for fine-tuning
		\end{itemize}
		
		\vspace{0.1cm}
		\textcolor{myred}{\textbf{WHY OTHER OPTIONS ARE WRONG:}}
		\begin{itemize}
			\item \textcolor{myred}{A:} Fine-tuning happens after pre-training
			\item \textcolor{myred}{C:} Deployment happens after training
			\item \textcolor{myred}{D:} Evaluation happens after training
		\end{itemize}
		
		\textcolor{myred}{\textbf{EXAM TRAP:}} Pre-training = \textcolor{myblue}{unlabeled text} learning!
	\end{frame}
	
	\begin{frame}{Q56: Fine-tuning LLMs - Tutorial}
		\fontsize{6.5}{7.5}\selectfont
		\textbf{QUESTION: What is the purpose of fine-tuning an LLM?}
		
		\vspace{0.1cm}
		\textbf{OPTIONS:}
		\begin{enumerate}
			\item To start training from scratch
			\item To transfer knowledge to specific tasks using labeled data
			\item To clean the data
			\item To tokenize the text
		\end{enumerate}
		
		\vspace{0.1cm}
		\textcolor{myblue}{\textbf{ANSWER: B}}
		
		\vspace{0.1cm}
		\textbf{DETAILED EXPLANATION:}
		\begin{itemize}
			\item Transfer knowledge from pre-training
			\item Use labeled data for specific tasks
			\item Adapt model to domain
			\item Enable task-specific performance
		\end{itemize}
		
		\vspace{0.1cm}
		\textcolor{myred}{\textbf{WHY OTHER OPTIONS ARE WRONG:}}
		\begin{itemize}
			\item \textcolor{myred}{A:} Fine-tuning uses existing pre-trained model
			\item \textcolor{myred}{C:} Data cleaning happens before training
			\item \textcolor{myred}{D:} Tokenization happens before training
		\end{itemize}
		
		\textcolor{myred}{\textbf{EXAM TRAP:}} Fine-tuning = \textcolor{myblue}{task-specific adaptation} with labeled data!
	\end{frame}
	
	\begin{frame}{Q57: RLHF Definition - Tutorial}
		\fontsize{6.5}{7.5}\selectfont
		\textbf{QUESTION: What is RLHF in LLM training?}
		
		\vspace{0.1cm}
		\textbf{OPTIONS:}
		\begin{enumerate}
			\item Reinforcement Learning Through Human Feedback
			\item Recurrent Learning with Human Feedback
			\item Regularization Learning with Human Feedback
			\item Regression Learning with Human Feedback
		\end{enumerate}
		
		\vspace{0.1cm}
		\textcolor{myblue}{\textbf{ANSWER: A}}
		
		\vspace{0.1cm}
		\textbf{DETAILED EXPLANATION:}
		\begin{itemize}
			\item Reinforcement Learning through Human Feedback
			\item Used in alignment phase
			\item Conforms model to human preferences
			\item Based on user feedback
		\end{itemize}
		
		\vspace{0.1cm}
		\textcolor{myred}{\textbf{WHY OTHER OPTIONS ARE WRONG:}}
		\begin{itemize}
			\item \textcolor{myred}{B:} Incorrect expansion of RLHF
			\item \textcolor{myred}{C:} Incorrect expansion of RLHF
			\item \textcolor{myred}{D:} Incorrect expansion of RLHF
		\end{itemize}
		
		\textcolor{myred}{\textbf{EXAM TRAP:}} RLHF = \textcolor{myblue}{Reinforcement Learning through Human Feedback}!
	\end{frame}
	
	\begin{frame}{Q58: Cloud-based LLMs - Tutorial}
		\fontsize{6.5}{7.5}\selectfont
		\textbf{QUESTION: Why are most LLMs cloud-based?}
		
		\vspace{0.1cm}
		\textbf{OPTIONS:}
		\begin{enumerate}
			\item They are too small to store locally
			\item They are huge (billions of parameters) and training is only done once
			\item It's cheaper to store locally
			\item They don't require storage
		\end{enumerate}
		
		\vspace{0.1cm}
		\textcolor{myblue}{\textbf{ANSWER: B}}
		
		\vspace{0.1cm}
		\textbf{DETAILED EXPLANATION:}
		\begin{itemize}
			\item Weight matrices are massive (billions of parameters)
			\item Infeasible to store locally
			\item Training only done once (too expensive to repeat)
			\item Reside remotely on cloud storage
		\end{itemize}
		
		\vspace{0.1cm}
		\textcolor{myred}{\textbf{WHY OTHER OPTIONS ARE WRONG:}}
		\begin{itemize}
			\item \textcolor{myred}{A:} They are huge, not small
			\item \textcolor{myred}{C:} Local storage is not feasible
			\item \textcolor{myred}{D:} They need storage (cloud storage)
		\end{itemize}
		
		\textcolor{myred}{\textbf{EXAM TRAP:}} LLMs are \textcolor{myblue}{cloud-based due to size}!
	\end{frame}
	
	\begin{frame}{Q59: GPT Evolution - Tutorial}
		\fontsize{6.5}{7.5}\selectfont
		\textbf{QUESTION: Which GPT generation introduced "no-shot fine-tuning"?}
		
		\vspace{0.1cm}
		\textbf{OPTIONS:}
		\begin{enumerate}
			\item GPT-1
			\item GPT-2
			\item GPT-3
			\item GPT-4
		\end{enumerate}
		
		\vspace{0.1cm}
		\textcolor{myblue}{\textbf{ANSWER: B}}
		
		\vspace{0.1cm}
		\textbf{DETAILED EXPLANATION:}
		\begin{itemize}
			\item GPT-2 removed need for second stage fine-tuning
			\item No-shot fine-tuning became possible
			\item Model adapts to new situations without examples
			\item Trained on larger, more varied corpus
		\end{itemize}
		
		\vspace{0.1cm}
		\textcolor{myred}{\textbf{WHY OTHER OPTIONS ARE WRONG:}}
		\begin{itemize}
			\item \textcolor{myred}{A:} GPT-1 required fine-tuning
			\item \textcolor{myred}{C:} GPT-3 improved but GPT-2 introduced no-shot
			\item \textcolor{myred}{D:} GPT-4 came later
		\end{itemize}
		
		\textcolor{myred}{\textbf{EXAM TRAP:}} GPT-2 introduced \textcolor{myblue}{no-shot fine-tuning}!
	\end{frame}
	
	\begin{frame}{Q60: No-shot Fine-tuning - Tutorial}
		\fontsize{6.5}{7.5}\selectfont
		\textbf{QUESTION: What is no-shot fine-tuning in LLMs?}
		
		\vspace{0.1cm}
		\textbf{OPTIONS:}
		\begin{enumerate}
			\item Requires many examples to adapt
			\item Adapts to new situations without any examples
			\item Requires exactly one example
			\item Cannot adapt to new situations
		\end{enumerate}
		
		\vspace{0.1cm}
		\textcolor{myblue}{\textbf{ANSWER: B}}
		
		\vspace{0.1cm}
		\textbf{DETAILED EXPLANATION:}
		\begin{itemize}
			\item No supervised examples needed
			\item Model adapts to new tasks
			\item Relies on pre-training knowledge
			\item GPT-2 enabled this
		\end{itemize}
		
		\vspace{0.1cm}
		\textcolor{myred}{\textbf{WHY OTHER OPTIONS ARE WRONG:}}
		\begin{itemize}
			\item \textcolor{myred}{A:} This is few-shot, not no-shot
			\item \textcolor{myred}{C:} This is one-shot, not no-shot
			\item \textcolor{myred}{D:} No-shot enables adaptation
		\end{itemize}
		
		\textcolor{myred}{\textbf{EXAM TRAP:}} No-shot = \textcolor{myblue}{no examples needed}!
	\end{frame}
	
	\begin{frame}{Q61: Chatbot Types - Tutorial}
		\fontsize{6.5}{7.5}\selectfont
		\textbf{QUESTION: Which type of chatbot is most advanced?}
		
		\vspace{0.1cm}
		\textbf{OPTIONS:}
		\begin{enumerate}
			\item FAQ-based chatbots
			\item Flow-based chatbots
			\item LLM-powered chatbots
			\item All are equally advanced
		\end{enumerate}
		
		\vspace{0.1cm}
		\textcolor{myblue}{\textbf{ANSWER: C}}
		
		\vspace{0.1cm}
		\textbf{DETAILED EXPLANATION:}
		\begin{itemize}
			\item FAQ-based: Standard responses, limited
			\item Flow-based: Interactive, multi-step
			\item LLM-powered: Deep understanding, human-like
			\item LLM-powered chatbots are the most advanced
		\end{itemize}
		
		\vspace{0.1cm}
		\textcolor{myred}{\textbf{WHY OTHER OPTIONS ARE WRONG:}}
		\begin{itemize}
			\item \textcolor{myred}{A:} FAQ-based are the simplest
			\item \textcolor{myred}{B:} Flow-based are intermediate
			\item \textcolor{myred}{D:} They have different levels of sophistication
		\end{itemize}
		
		\textcolor{myred}{\textbf{EXAM TRAP:}} LLM-powered chatbots are \textcolor{myblue}{most advanced}!
	\end{frame}
	
	\begin{frame}{Q62: Prompt Engineering Definition - Tutorial}
		\fontsize{6.5}{7.5}\selectfont
		\textbf{QUESTION: What is prompt engineering?}
		
		\vspace{0.1cm}
		\textbf{OPTIONS:}
		\begin{enumerate}
			\item Training the model from scratch
			\item The art and science of designing prompts for effective LLM use
			\item Evaluating model performance
			\item Tokenizing text
		\end{enumerate}
		
		\vspace{0.1cm}
		\textcolor{myblue}{\textbf{ANSWER: B}}
		
		\vspace{0.1cm}
		\textbf{DETAILED EXPLANATION:}
		\begin{itemize}
			\item Designing prompts for best results
			\item Iterative process
			\item Provide context and examples
			\item Guide model responses
		\end{itemize}
		
		\vspace{0.1cm}
		\textcolor{myred}{\textbf{WHY OTHER OPTIONS ARE WRONG:}}
		\begin{itemize}
			\item \textcolor{myred}{A:} Training is different from prompting
			\item \textcolor{myred}{C:} Evaluation is different from prompting
			\item \textcolor{myred}{D:} Tokenization is a pre-processing step
		\end{itemize}
		
		\textcolor{myred}{\textbf{EXAM TRAP:}} Prompt engineering = \textcolor{myblue}{designing effective prompts}!
	\end{frame}
	
	\begin{frame}{Q63: Zero-shot Prompting - Tutorial}
		\fontsize{6.5}{7.5}\selectfont
		\textbf{QUESTION: What is zero-shot prompting?}
		
		\vspace{0.1cm}
		\textbf{OPTIONS:}
		\begin{enumerate}
			\item Providing many examples before the question
			\item Providing no examples before the question
			\item Providing one example before the question
			\item Providing multiple examples with reasoning
		\end{enumerate}
		
		\vspace{0.1cm}
		\textcolor{myblue}{\textbf{ANSWER: B}}
		
		\vspace{0.1cm}
		\textbf{DETAILED EXPLANATION:}
		\begin{itemize}
			\item No examples provided
			\item User inputs question directly
			\item Most common type
			\item "What is the capital of France?"
		\end{itemize}
		
		\vspace{0.1cm}
		\textcolor{myred}{\textbf{WHY OTHER OPTIONS ARE WRONG:}}
		\begin{itemize}
			\item \textcolor{myred}{A:} This is few-shot prompting
			\item \textcolor{myred}{C:} This is one-shot prompting
			\item \textcolor{myred}{D:} This is few-shot CoT
		\end{itemize}
		
		\textcolor{myred}{\textbf{EXAM TRAP:}} Zero-shot = \textcolor{myblue}{no examples}!
	\end{frame}
	



\begin{frame}{Q64: Few-shot Prompting - Tutorial}
	\fontsize{6.5}{7.5}\selectfont
	
	\textbf{QUESTION: What is few-shot prompting?}
	
	\vspace{0.1cm}
	\textbf{OPTIONS:}
	\begin{enumerate}
		\item Providing no examples
		\item Providing one example
		\item Providing multiple examples before the target question
		\item Providing examples after the answer
	\end{enumerate}
	
	\vspace{0.1cm}
	\textcolor{myblue}{\textbf{ANSWER: C}}
	
	\vspace{0.1cm}
	\textbf{DETAILED EXPLANATION:}
	\begin{itemize}
		\item Provides multiple examples in the prompt
		\item Gives the model more context
		\item Helps produce more accurate and consistent responses
		\item Example: Several Q\&A pairs are provided before the target question
	\end{itemize}
	
	\vspace{0.1cm}
	\textcolor{myred}{\textbf{WHY OTHER OPTIONS ARE WRONG:}}
	\begin{itemize}
		\item \textcolor{myred}{A:} This is zero-shot prompting
		\item \textcolor{myred}{B:} This is one-shot prompting
		\item \textcolor{myred}{D:} Examples are provided before, not after, the target question
	\end{itemize}
	
	\textcolor{myred}{\textbf{EXAM TRAP:}}
	Few-shot = \textcolor{myblue}{multiple examples before the target}!
\end{frame}



	
	\begin{frame}{Q65: Chain of Thought Prompting - Tutorial}
		\fontsize{6.5}{7.5}\selectfont
		\textbf{QUESTION: What is Chain of Thought (CoT) prompting?}
		
		\vspace{0.1cm}
		\textbf{OPTIONS:}
		\begin{enumerate}
			\item Prompting with the final answer only
			\item Prompting with intermediate reasoning steps
			\item Prompting with no text
			\item Prompting with random words
		\end{enumerate}
		
		\vspace{0.1cm}
		\textcolor{myblue}{\textbf{ANSWER: B}}
		
		\vspace{0.1cm}
		\textbf{DETAILED EXPLANATION:}
		\begin{itemize}
			\item Include reasoning steps in prompt
			\item Solves complex reasoning tasks
			\item Improves multi-step reasoning
			\item "Let's think step by step"
		\end{itemize}
		
		\vspace{0.1cm}
		\textcolor{myred}{\textbf{WHY OTHER OPTIONS ARE WRONG:}}
		\begin{itemize}
			\item \textcolor{myred}{A:} Standard prompting gives final answer only
			\item \textcolor{myred}{C:} CoT requires text
			\item \textcolor{myred}{D:} CoT uses structured reasoning
		\end{itemize}
		
		\textcolor{myred}{\textbf{EXAM TRAP:}} CoT = \textcolor{myblue}{intermediate reasoning steps}!
	\end{frame}
	
	\begin{frame}{Q66: Context Length - Tutorial}
		\fontsize{6.5}{7.5}\selectfont
		\textbf{QUESTION: What is context length in LLMs?}
		
		\vspace{0.1cm}
		\textbf{OPTIONS:}
		\begin{enumerate}
			\item The number of parameters in the model
			\item The maximum number of tokens a model can process in a single sequence
			\item The number of layers in the model
			\item The training time of the model
		\end{enumerate}
		
		\vspace{0.1cm}
		\textcolor{myblue}{\textbf{ANSWER: B}}
		
		\vspace{0.1cm}
		\textbf{DETAILED EXPLANATION:}
		\begin{itemize}
			\item Maximum tokens in input/output sequence
			\item Also called context window
			\item Limits model's working memory
			\item Important for long documents
		\end{itemize}
		
		\vspace{0.1cm}
		\textcolor{myred}{\textbf{WHY OTHER OPTIONS ARE WRONG:}}
		\begin{itemize}
			\item \textcolor{myred}{A:} This is parameter count
			\item \textcolor{myred}{C:} This is model depth
			\item \textcolor{myred}{D:} This is training duration
		\end{itemize}
		
		\textcolor{myred}{\textbf{EXAM TRAP:}} Context length = \textcolor{myblue}{tokens model can process}!
	\end{frame}
	
	\begin{frame}{Q67: Statelessness Definition - Tutorial}
		\fontsize{6.5}{7.5}\selectfont
		\textbf{QUESTION: What does statelessness mean in LLMs?}
		
		\vspace{0.1cm}
		\textbf{OPTIONS:}
		\begin{enumerate}
			\item LLMs remember all past conversations
			\item LLMs don't retain information between prompts
			\item LLMs retain information for one day
			\item LLMs have perfect memory
		\end{enumerate}
		
		\vspace{0.1cm}
		\textcolor{myblue}{\textbf{ANSWER: B}}
		
		\vspace{0.1cm}
		\textbf{DETAILED EXPLANATION:}
		\begin{itemize}
			\item Each prompt processed independently
			\item No memory of previous conversations
			\item By design
			\item Users must feed previous responses for context
		\end{itemize}
		
		\vspace{0.1cm}
		\textcolor{myred}{\textbf{WHY OTHER OPTIONS ARE WRONG:}}
		\begin{itemize}
			\item \textcolor{myred}{A:} This is stateful (opposite)
			\item \textcolor{myred}{C:} There is no time-based retention
			\item \textcolor{myred}{D:} LLMs have no perfect memory
		\end{itemize}
		
		\textcolor{myred}{\textbf{EXAM TRAP:}} Stateless = \textcolor{myblue}{no memory between prompts}!
	\end{frame}
	
	\begin{frame}{Q68: Temperature Parameter - Tutorial}
		\fontsize{6.5}{7.5}\selectfont
		\textbf{QUESTION: What does the temperature parameter control in LLMs?}
		
		\vspace{0.1cm}
		\textbf{OPTIONS:}
		\begin{enumerate}
			\item The number of tokens generated
			\item The randomness/creativity of output
			\item The length of the output
			\item The speed of generation
		\end{enumerate}
		
		\vspace{0.1cm}
		\textcolor{myblue}{\textbf{ANSWER: B}}
		
		\vspace{0.1cm}
		\textbf{DETAILED EXPLANATION:}
		\begin{itemize}
			\item Controls randomness/creativity
			\item Shapes probability distribution
			\item Lower = more predictable, deterministic
			\item Higher = more creative, varied
		\end{itemize}
		
		\vspace{0.1cm}
		\textcolor{myred}{\textbf{WHY OTHER OPTIONS ARE WRONG:}}
		\begin{itemize}
			\item \textcolor{myred}{A:} Token count is a separate parameter
			\item \textcolor{myred}{C:} Output length is a separate parameter
			\item \textcolor{myred}{D:} Speed is determined by hardware
		\end{itemize}
		
		\textcolor{myred}{\textbf{EXAM TRAP:}} Temperature controls \textcolor{myblue}{randomness/creativity}!
	\end{frame}
	
	\begin{frame}{Q69: Top-K Sampling - Tutorial}
		\fontsize{6.5}{7.5}\selectfont
		\textbf{QUESTION: What does Top-K sampling do in LLMs?}
		
		\vspace{0.1cm}
		\textbf{OPTIONS:}
		\begin{enumerate}
			\item Considers all tokens equally
			\item Only considers the K most probable tokens
			\item Randomly selects tokens
			\item Always selects the most probable token
		\end{enumerate}
		
		\vspace{0.1cm}
		\textcolor{myblue}{\textbf{ANSWER: B}}
		
		\vspace{0.1cm}
		\textbf{DETAILED EXPLANATION:}
		\begin{itemize}
			\item Restricts to K most probable tokens
			\item Low K = less creativity
			\item High K = more variety
			\item Example: K=3 means only top 3 considered
		\end{itemize}
		
		\vspace{0.1cm}
		\textcolor{myred}{\textbf{WHY OTHER OPTIONS ARE WRONG:}}
		\begin{itemize}
			\item \textcolor{myred}{A:} Top-K restricts, doesn't consider all
			\item \textcolor{myred}{C:} Top-K is deterministic, not random
			\item \textcolor{myred}{D:} This is temperature=0, not Top-K
		\end{itemize}
		
		\textcolor{myred}{\textbf{EXAM TRAP:}} Top-K = \textcolor{myblue}{only K most probable tokens}!
	\end{frame}
	
	\begin{frame}{Q70: Agentic AI Definition - Tutorial}
		\fontsize{6.5}{7.5}\selectfont
		\textbf{QUESTION: What is Agentic AI?}
		
		\vspace{0.1cm}
		\textbf{OPTIONS:}
		\begin{enumerate}
			\item AI that only generates text
			\item AI agents that take inputs and perform actions autonomously
			\item AI that only classifies data
			\item AI that requires constant human supervision
		\end{enumerate}
		
		\vspace{0.1cm}
		\textcolor{myblue}{\textbf{ANSWER: B}}
		
		\vspace{0.1cm}
		\textbf{DETAILED EXPLANATION:}
		\begin{itemize}
			\item Next frontier in AI
			\item Autonomous action
			\item Uses LLM outputs
			\item Dynamic decision-making
		\end{itemize}
		
		\vspace{0.1cm}
		\textcolor{myred}{\textbf{WHY OTHER OPTIONS ARE WRONG:}}
		\begin{itemize}
			\item \textcolor{myred}{A:} Agentic AI goes beyond text generation
			\item \textcolor{myred}{C:} Agentic AI takes action, not just classify
			\item \textcolor{myred}{D:} Agentic AI operates with limited supervision
		\end{itemize}
		
		\textcolor{myred}{\textbf{EXAM TRAP:}} Agentic AI = \textcolor{myblue}{autonomous action}!
	\end{frame}
	
	% ============================================================
	% SECTION 6: GENAI PERFORMANCE ANALYSIS - QUESTIONS 71-85
	% ============================================================
	
	\begin{frame}{Q71: GenAI Evaluation Challenge - Tutorial}
		\fontsize{6.5}{7.5}\selectfont
		\textbf{QUESTION: Why is GenAI evaluation more complex than traditional ML?}
		
		\vspace{0.1cm}
		\textbf{OPTIONS:}
		\begin{enumerate}
			\item GenAI outputs are always correct
			\item GenAI outputs are open-ended and unstructured
			\item GenAI outputs are always deterministic
			\item GenAI outputs are simple to evaluate
		\end{enumerate}
		
		\vspace{0.1cm}
		\textcolor{myblue}{\textbf{ANSWER: B}}
		
		\vspace{0.1cm}
		\textbf{DETAILED EXPLANATION:}
		\begin{itemize}
			\item Unstructured responses
			\item No single correct answer
			\item Many valid variations
			\item Hallucination risk
		\end{itemize}
		
		\vspace{0.1cm}
		\textcolor{myred}{\textbf{WHY OTHER OPTIONS ARE WRONG:}}
		\begin{itemize}
			\item \textcolor{myred}{A:} GenAI outputs can be incorrect
			\item \textcolor{myred}{C:} GenAI outputs are probabilistic
			\item \textcolor{myred}{D:} Evaluation is more complex, not simple
		\end{itemize}
		
		\textcolor{myred}{\textbf{EXAM TRAP:}} GenAI evaluation is \textcolor{myblue}{complex due to open-ended outputs}!
	\end{frame}
	
	\begin{frame}{Q72: Hallucination Definition - Tutorial}
		\fontsize{6.5}{7.5}\selectfont
		\textbf{QUESTION: What is hallucination in GenAI?}
		
		\vspace{0.1cm}
		\textbf{OPTIONS:}
		\begin{enumerate}
			\item Output that is always correct
			\item Output that appears reasonable but is factually incorrect
			\item Output that is completely nonsensical
			\item Output that is copied verbatim from training data
		\end{enumerate}
		
		\vspace{0.1cm}
		\textcolor{myblue}{\textbf{ANSWER: B}}
		
		\vspace{0.1cm}
		\textbf{DETAILED EXPLANATION:}
		\begin{itemize}
			\item Plausible but incorrect
			\item Not supported by input
			\item Caused by probabilistic generation
			\item A serious GenAI risk
		\end{itemize}
		
		\vspace{0.1cm}
		\textcolor{myred}{\textbf{WHY OTHER OPTIONS ARE WRONG:}}
		\begin{itemize}
			\item \textcolor{myred}{A:} Hallucinations are factually incorrect
			\item \textcolor{myred}{C:} Hallucinations appear reasonable, not nonsensical
			\item \textcolor{myred}{D:} Hallucinations are not necessarily copied
		\end{itemize}
		
		\textcolor{myred}{\textbf{EXAM TRAP:}} Hallucination = \textcolor{myblue}{reasonable but incorrect} output!
	\end{frame}
	
	\begin{frame}{Q73: Hallucination Causes - Tutorial}
		\fontsize{6.5}{7.5}\selectfont
		\textbf{QUESTION: What is a primary cause of hallucination in LLMs?}
		
		\vspace{0.1cm}
		\textbf{OPTIONS:}
		\begin{enumerate}
			\item Perfect training data
			\item Probabilistic token selection
			\item Deterministic outputs
			\item Small model size
		\end{enumerate}
		
		\vspace{0.1cm}
		\textcolor{myblue}{\textbf{ANSWER: B}}
		
		\vspace{0.1cm}
		\textbf{DETAILED EXPLANATION:}
		\begin{itemize}
			\item Tokens chosen based on probabilities
			\item Higher temperature = more variation
			\item May pick incorrect tokens
			\item Knowledge gaps filled with plausible answers
		\end{itemize}
		
		\vspace{0.1cm}
		\textcolor{myred}{\textbf{WHY OTHER OPTIONS ARE WRONG:}}
		\begin{itemize}
			\item \textcolor{myred}{A:} Training data is never perfect
			\item \textcolor{myred}{C:} LLM outputs are probabilistic, not deterministic
			\item \textcolor{myred}{D:} Larger models can also hallucinate
		\end{itemize}
		
		\textcolor{myred}{\textbf{EXAM TRAP:}} Hallucination = \textcolor{myblue}{probabilistic nature} + knowledge gaps!
	\end{frame}
	
	\begin{frame}{Q74: Robustness Evaluation - Tutorial}
		\fontsize{6.5}{7.5}\selectfont
		\textbf{QUESTION: What does robustness in GenAI evaluation measure?}
		
		\vspace{0.1cm}
		\textbf{OPTIONS:}
		\begin{enumerate}
			\item How fast the model runs
			\item How well the model performs on novel/unseen inputs
			\item How many parameters the model has
			\item How much data the model was trained on
		\end{enumerate}
		
		\vspace{0.1cm}
		\textcolor{myblue}{\textbf{ANSWER: B}}
		
		\vspace{0.1cm}
		\textbf{DETAILED EXPLANATION:}
		\begin{itemize}
			\item Out of Distribution (OOD) performance
			\item Adversarial robustness
			\item Jailbreak resistance
			\item Prompt sensitivity
		\end{itemize}
		
		\vspace{0.1cm}
		\textcolor{myred}{\textbf{WHY OTHER OPTIONS ARE WRONG:}}
		\begin{itemize}
			\item \textcolor{myred}{A:} This is operational performance
			\item \textcolor{myred}{C:} This is model size, not robustness
			\item \textcolor{myred}{D:} This is training data size
		\end{itemize}
		
		\textcolor{myred}{\textbf{EXAM TRAP:}} Robustness = \textcolor{myblue}{performance on unseen/attacked inputs}!
	\end{frame}
	
	\begin{frame}{Q75: Bias Metrics - Tutorial}
		\fontsize{6.5}{7.5}\selectfont
		\textbf{QUESTION: What is WEAT used to measure?}
		
		\vspace{0.1cm}
		\textbf{OPTIONS:}
		\begin{enumerate}
			\item Model accuracy
			\item Bias using cosine similarity between sets of words
			\item Model speed
			\item Model size
		\end{enumerate}
		
		\vspace{0.1cm}
		\textcolor{myblue}{\textbf{ANSWER: B}}
		
		\vspace{0.1cm}
		\textbf{DETAILED EXPLANATION:}
		\begin{itemize}
			\item WEAT = Word Embedding Association Test
			\item Measures bias using cosine similarity
			\item Between sets of words
			\item Normalized score
		\end{itemize}
		
		\vspace{0.1cm}
		\textcolor{myred}{\textbf{WHY OTHER OPTIONS ARE WRONG:}}
		\begin{itemize}
			\item \textcolor{myred}{A:} WEAT measures bias, not accuracy
			\item \textcolor{myred}{C:} WEAT doesn't measure speed
			\item \textcolor{myred}{D:} WEAT doesn't measure size
		\end{itemize}
		
		\textcolor{myred}{\textbf{EXAM TRAP:}} WEAT = \textcolor{myblue}{bias measurement} with cosine similarity!
	\end{frame}
	
	\begin{frame}{Q76: Operational Performance Metrics - Tutorial}
		\fontsize{6.5}{7.5}\selectfont
		\textbf{QUESTION: What is "Time to First Token" in LLM evaluation?}
		
		\vspace{0.1cm}
		\textbf{OPTIONS:}
		\begin{enumerate}
			\item Total generation time
			\item Time between request and first token generation
			\item Time between tokens
			\item Time to complete the entire response
		\end{enumerate}
		
		\vspace{0.1cm}
		\textcolor{myblue}{\textbf{ANSWER: B}}
		
		\vspace{0.1cm}
		\textbf{DETAILED EXPLANATION:}
		\begin{itemize}
			\item Perceived responsiveness
			\item Important for user experience
			\item Measures initial delay
			\item Affects user satisfaction
		\end{itemize}
		
		\vspace{0.1cm}
		\textcolor{myred}{\textbf{WHY OTHER OPTIONS ARE WRONG:}}
		\begin{itemize}
			\item \textcolor{myred}{A:} Total generation time is different
			\item \textcolor{myred}{C:} This is time between tokens
			\item \textcolor{myred}{D:} This is total time, not first token
		\end{itemize}
		
		\textcolor{myred}{\textbf{EXAM TRAP:}} Time to First Token = \textcolor{myblue}{initial response delay}!
	\end{frame}
	
	\begin{frame}{Q77: BLEU Score - Tutorial}
		\fontsize{6.5}{7.5}\selectfont
		\textbf{QUESTION: What does BLEU score measure in text generation?}
		
		\vspace{0.1cm}
		\textbf{OPTIONS:}
		\begin{enumerate}
			\item Creativity of output
			\item n-gram precision compared to reference text
			\item Speed of generation
			\item Model size
		\end{enumerate}
		
		\vspace{0.1cm}
		\textcolor{myblue}{\textbf{ANSWER: B}}
		
		\vspace{0.1cm}
		\textbf{DETAILED EXPLANATION:}
		\begin{itemize}
			\item BLEU = Bilingual Evaluation Understudy
			\item n-gram precision
			\item Adds brevity penalty
			\item 0-1 scale
		\end{itemize}
		
		\vspace{0.1cm}
		\textcolor{myred}{\textbf{WHY OTHER OPTIONS ARE WRONG:}}
		\begin{itemize}
			\item \textcolor{myred}{A:} BLEU measures similarity, not creativity
			\item \textcolor{myred}{C:} BLEU doesn't measure speed
			\item \textcolor{myred}{D:} BLEU doesn't measure model size
		\end{itemize}
		
		\textcolor{myred}{\textbf{EXAM TRAP:}} BLEU = \textcolor{myblue}{n-gram precision} vs reference!
	\end{frame}
	
	\begin{frame}{Q78: ROUGE Score - Tutorial}
		\fontsize{6.5}{7.5}\selectfont
		\textbf{QUESTION: What is ROUGE used for in GenAI evaluation?}
		
		\vspace{0.1cm}
		\textbf{OPTIONS:}
		\begin{enumerate}
			\item Evaluating image quality
			\item Evaluating summaries/translations against reference
			\item Evaluating audio quality
			\item Evaluating model speed
		\end{enumerate}
		
		\vspace{0.1cm}
		\textcolor{myblue}{\textbf{ANSWER: B}}
		
		\vspace{0.1cm}
		\textbf{DETAILED EXPLANATION:}
		\begin{itemize}
			\item ROUGE = Recall-Oriented Understudy
			\item For summaries and translations
			\item n-gram matching
			\item 0-1 scale
		\end{itemize}
		
		\vspace{0.1cm}
		\textcolor{myred}{\textbf{WHY OTHER OPTIONS ARE WRONG:}}
		\begin{itemize}
			\item \textcolor{myred}{A:} ROUGE is for text, not images
			\item \textcolor{myred}{C:} ROUGE is for text, not audio
			\item \textcolor{myred}{D:} ROUGE measures quality, not speed
		\end{itemize}
		
		\textcolor{myred}{\textbf{EXAM TRAP:}} ROUGE = \textcolor{myblue}{summary/translation evaluation}!
	\end{frame}
	
	\begin{frame}{Q79: FID Score - Tutorial}
		\fontsize{6.5}{7.5}\selectfont
		\textbf{QUESTION: What does FID measure in image generation?}
		
		\vspace{0.1cm}
		\textbf{OPTIONS:}
		\begin{enumerate}
			\item How fast images are generated
			\item How close generated images are to real images
			\item How large images are
			\item How many colors images have
		\end{enumerate}
		
		\vspace{0.1cm}
		\textcolor{myblue}{\textbf{ANSWER: B}}
		
		\vspace{0.1cm}
		\textbf{DETAILED EXPLANATION:}
		\begin{itemize}
			\item FID = Fréchet Inception Distance
			\item Compares generated vs real images
			\item Distribution parameters (mean, covariance)
			\item Lower FID = better quality
		\end{itemize}
		
		\vspace{0.1cm}
		\textcolor{myred}{\textbf{WHY OTHER OPTIONS ARE WRONG:}}
		\begin{itemize}
			\item \textcolor{myred}{A:} FID measures quality, not speed
			\item \textcolor{myred}{C:} FID doesn't measure image size
			\item \textcolor{myred}{D:} FID doesn't measure color count
		\end{itemize}
		
		\textcolor{myred}{\textbf{EXAM TRAP:}} FID = \textcolor{myblue}{generated vs real image comparison}!
	\end{frame}
	
	\begin{frame}{Q80: MMLU Benchmark - Tutorial}
		\fontsize{6.5}{7.5}\selectfont
		\textbf{QUESTION: What does MMLU benchmark evaluate?}
		
		\vspace{0.1cm}
		\textbf{OPTIONS:}
		\begin{enumerate}
			\item Image generation quality
			\item Language understanding across 57 subjects
			\item Audio quality
			\item Model speed
		\end{enumerate}
		
		\vspace{0.1cm}
		\textcolor{myblue}{\textbf{ANSWER: B}}
		
		\vspace{0.1cm}
		\textbf{DETAILED EXPLANATION:}
		\begin{itemize}
			\item MMLU = Massive Multitask Language Understanding
			\item Evaluates across 57 subjects
			\item Humanities, social sciences, STEM
			\item Multiple-choice format
		\end{itemize}
		
		\vspace{0.1cm}
		\textcolor{myred}{\textbf{WHY OTHER OPTIONS ARE WRONG:}}
		\begin{itemize}
			\item \textcolor{myred}{A:} MMLU is text-based, not image
			\item \textcolor{myred}{C:} MMLU is text-based, not audio
			\item \textcolor{myred}{D:} MMLU measures knowledge, not speed
		\end{itemize}
		
		\textcolor{myred}{\textbf{EXAM TRAP:}} MMLU = \textcolor{myblue}{multitask language understanding} across 57 subjects!
	\end{frame}
	
	\begin{frame}{Q81: Data Leakage in Benchmarks - Tutorial}
		\fontsize{6.5}{7.5}\selectfont
		\textbf{QUESTION: What is data leakage in LLM benchmarking?}
		
		\vspace{0.1cm}
		\textbf{OPTIONS:}
		\begin{enumerate}
			\item Benchmark data being included in training data
			\item Perfect model performance
			\item Model not being tested
			\item Model being too slow
		\end{enumerate}
		
		\vspace{0.1cm}
		\textcolor{myblue}{\textbf{ANSWER: A}}
		
		\vspace{0.1cm}
		\textbf{DETAILED EXPLANATION:}
		\begin{itemize}
			\item Benchmark data in training set
			\item Model "memorizes" answers
			\item Inflated performance
			\item False indication of ability
		\end{itemize}
		
		\vspace{0.1cm}
		\textcolor{myred}{\textbf{WHY OTHER OPTIONS ARE WRONG:}}
		\begin{itemize}
			\item \textcolor{myred}{B:} Data leakage leads to false performance
			\item \textcolor{myred}{C:} Data leakage means model IS tested (incorrectly)
			\item \textcolor{myred}{D:} Data leakage is about data contamination, not speed
		\end{itemize}
		
		\textcolor{myred}{\textbf{EXAM TRAP:}} Data leakage = \textcolor{myblue}{benchmark data in training}!
	\end{frame}
	
	\begin{frame}{Q82: Benchmark Saturation - Tutorial}
		\fontsize{6.5}{7.5}\selectfont
		\textbf{QUESTION: What is benchmark saturation?}
		
		\vspace{0.1cm}
		\textbf{OPTIONS:}
		\begin{enumerate}
			\item Models cannot solve the benchmark
			\item Models reach maximum scores, making differentiation impossible
			\item The benchmark is too difficult
			\item The benchmark is not used
		\end{enumerate}
		
		\vspace{0.1cm}
		\textcolor{myblue}{\textbf{ANSWER: B}}
		
		\vspace{0.1cm}
		\textbf{DETAILED EXPLANATION:}
		\begin{itemize}
			\item Models reach ceiling scores
			\item Cannot differentiate performance
			\item Benchmarks become outdated
			\item Need new challenges
		\end{itemize}
		
		\vspace{0.1cm}
		\textcolor{myred}{\textbf{WHY OTHER OPTIONS ARE WRONG:}}
		\begin{itemize}
			\item \textcolor{myred}{A:} Saturation means models CAN solve it
			\item \textcolor{myred}{C:} Saturation means it's too easy, not too difficult
			\item \textcolor{myred}{D:} Saturation means it's used but ineffective
		\end{itemize}
		
		\textcolor{myred}{\textbf{EXAM TRAP:}} Saturation = \textcolor{myblue}{models max out benchmark}!
	\end{frame}
	
	\begin{frame}{Q83: Cherry Picking in Benchmarks - Tutorial}
		\fontsize{6.5}{7.5}\selectfont
		\textbf{QUESTION: What is cherry picking in LLM benchmark reporting?}
		
		\vspace{0.1cm}
		\textbf{OPTIONS:}
		\begin{enumerate}
			\item Reporting all benchmark results fairly
			\item Only highlighting good benchmark results, omitting weaker ones
			\item Using the latest benchmarks
			\item Comparing models equally
		\end{enumerate}
		
		\vspace{0.1cm}
		\textcolor{myblue}{\textbf{ANSWER: B}}
		
		\vspace{0.1cm}
		\textbf{DETAILED EXPLANATION:}
		\begin{itemize}
			\item Selectively report good results
			\item Omit weaker benchmarks
			\item Misleading representation
			\item Selective reporting
		\end{itemize}
		
		\vspace{0.1cm}
		\textcolor{myred}{\textbf{WHY OTHER OPTIONS ARE WRONG:}}
		\begin{itemize}
			\item \textcolor{myred}{A:} This is fair reporting, not cherry picking
			\item \textcolor{myred}{C:} Cherry picking is about selection, not recency
			\item \textcolor{myred}{D:} Cherry picking is unfair comparison
		\end{itemize}
		
		\textcolor{myred}{\textbf{EXAM TRAP:}} Cherry picking = \textcolor{myblue}{selective reporting} of good results!
	\end{frame}
	
	\begin{frame}{Q84: Hallucination Rate - Tutorial}
		\fontsize{6.5}{7.5}\selectfont
		\textbf{QUESTION: What does Hallucination Rate measure?}
		
		\vspace{0.1cm}
		\textbf{OPTIONS:}
		\begin{enumerate}
			\item Model speed
			\item Model's propensity for inaccurate output
			\item Model size
			\item Model training time
		\end{enumerate}
		
		\vspace{0.1cm}
		\textcolor{myblue}{\textbf{ANSWER: B}}
		
		\vspace{0.1cm}
		\textbf{DETAILED EXPLANATION:}
		\begin{itemize}
			\item Quantifies hallucination frequency
			\item Essential risk metric
			\item Identifies unreliable outputs
			\item Critical for risk management
		\end{itemize}
		
		\vspace{0.1cm}
		\textcolor{myred}{\textbf{WHY OTHER OPTIONS ARE WRONG:}}
		\begin{itemize}
			\item \textcolor{myred}{A:} Hallucination rate is about quality, not speed
			\item \textcolor{myred}{C:} Hallucination rate is about accuracy, not size
			\item \textcolor{myred}{D:} Hallucination rate is about output quality, not training
		\end{itemize}
		
		\textcolor{myred}{\textbf{EXAM TRAP:}} Hallucination Rate = \textcolor{myblue}{propensity for inaccurate output}!
	\end{frame}
	
	\begin{frame}{Q85: Operational Performance Summary - Tutorial}
		\fontsize{6.5}{7.5}\selectfont
		\textbf{QUESTION: What are key operational performance metrics for LLMs?}
		
		\vspace{0.1cm}
		\textbf{OPTIONS:}
		\begin{enumerate}
			\item Accuracy and precision only
			\item Latency, throughput, and token efficiency
			\item Model size only
			\item Training data size only
		\end{enumerate}
		
		\vspace{0.1cm}
		\textcolor{myblue}{\textbf{ANSWER: B}}
		
		\vspace{0.1cm}
		\textbf{DETAILED EXPLANATION:}
		\begin{itemize}
			\item \textbf{Latency:} Response time
			\item \textbf{Throughput:} Queries per time
			\item \textbf{Token efficiency:} Quality per token
			\item Critical for production deployment
		\end{itemize}
		
		\vspace{0.1cm}
		\textcolor{myred}{\textbf{WHY OTHER OPTIONS ARE WRONG:}}
		\begin{itemize}
			\item \textcolor{myred}{A:} Accuracy/precision are quality metrics
			\item \textcolor{myred}{C:} Model size is a characteristic, not performance
			\item \textcolor{myred}{D:} Training data size is a characteristic
		\end{itemize}
		
		\textcolor{myred}{\textbf{EXAM TRAP:}} Operational metrics = \textcolor{myblue}{latency, throughput, token efficiency}!
	\end{frame}
	
	% ============================================================
	% SECTION 7: APPLICATIONS - QUESTIONS 86-100
	% ============================================================
	
	\begin{frame}{Q86: GenAI Text Applications - Tutorial}
		\fontsize{6.5}{7.5}\selectfont
		\textbf{QUESTION: What is a common text processing application of GenAI?}
		
		\vspace{0.1cm}
		\textbf{OPTIONS:}
		\begin{enumerate}
			\item Image generation
			\item Document summarization
			\item Video creation
			\item Audio synthesis
		\end{enumerate}
		
		\vspace{0.1cm}
		\textcolor{myblue}{\textbf{ANSWER: B}}
		
		\vspace{0.1cm}
		\textbf{DETAILED EXPLANATION:}
		\begin{itemize}
			\item Report generation
			\item Document summarization
			\item Legal analysis support
			\item Writing assistance
		\end{itemize}
		
		\vspace{0.1cm}
		\textcolor{myred}{\textbf{WHY OTHER OPTIONS ARE WRONG:}}
		\begin{itemize}
			\item \textcolor{myred}{A:} Image generation is a different modality
			\item \textcolor{myred}{C:} Video creation is a different modality
			\item \textcolor{myred}{D:} Audio synthesis is a different modality
		\end{itemize}
		
		\textcolor{myred}{\textbf{EXAM TRAP:}} Text processing includes \textcolor{myblue}{summarization and report generation}!
	\end{frame}
	
	\begin{frame}{Q87: GenAI in Compliance - Tutorial}
		\fontsize{6.5}{7.5}\selectfont
		\textbf{QUESTION: How can GenAI help with compliance monitoring?}
		
		\vspace{0.1cm}
		\textbf{OPTIONS:}
		\begin{enumerate}
			\item By ignoring all regulations
			\item By analyzing SEC filings and detecting emerging risks
			\item By removing all compliance requirements
			\item By creating new regulations
		\end{enumerate}
		
		\vspace{0.1cm}
		\textcolor{myblue}{\textbf{ANSWER: B}}
		
		\vspace{0.1cm}
		\textbf{DETAILED EXPLANATION:}
		\begin{itemize}
			\item Analyze SEC filings
			\item Review company reports
			\item Detect emerging risks
			\item Regulatory compliance
		\end{itemize}
		
		\vspace{0.1cm}
		\textcolor{myred}{\textbf{WHY OTHER OPTIONS ARE WRONG:}}
		\begin{itemize}
			\item \textcolor{myred}{A:} Compliance requires attention to regulations
			\item \textcolor{myred}{C:} Requirements cannot be ignored
			\item \textcolor{myred}{D:} GenAI doesn't create regulations
		\end{itemize}
		
		\textcolor{myred}{\textbf{EXAM TRAP:}} GenAI helps with \textcolor{myblue}{compliance monitoring}!
	\end{frame}
	
	\begin{frame}{Q88: Specialized Finance LLMs - Tutorial}
		\fontsize{6.5}{7.5}\selectfont
		\textbf{QUESTION: What are examples of finance-specialized LLMs?}
		
		\vspace{0.1cm}
		\textbf{OPTIONS:}
		\begin{enumerate}
			\item ChatGPT and Claude
			\item Bloomberg GPT and FinGPT
			\item DALL-E and Stable Diffusion
			\item Sora and Gemini
		\end{enumerate}
		
		\vspace{0.1cm}
		\textcolor{myblue}{\textbf{ANSWER: B}}
		
		\vspace{0.1cm}
		\textbf{DETAILED EXPLANATION:}
		\begin{itemize}
			\item Bloomberg GPT
			\item FinGPT
			\item Financial news analysis
			\item Trading signal generation
		\end{itemize}
		
		\vspace{0.1cm}
		\textcolor{myred}{\textbf{WHY OTHER OPTIONS ARE WRONG:}}
		\begin{itemize}
			\item \textcolor{myred}{A:} These are general-purpose LLMs
			\item \textcolor{myred}{C:} These are image generation models
			\item \textcolor{myred}{D:} These are multimodal/general models
		\end{itemize}
		
		\textcolor{myred}{\textbf{EXAM TRAP:}} Finance-specific LLMs = \textcolor{myblue}{Bloomberg GPT, FinGPT}!
	\end{frame}
	
	\begin{frame}{Q89: Code Generation Models - Tutorial}
		\fontsize{6.5}{7.5}\selectfont
		\textbf{QUESTION: Which model is specialized for code generation?}
		
		\vspace{0.1cm}
		\textbf{OPTIONS:}
		\begin{enumerate}
			\item Codex (OpenAI)
			\item DALL-E
			\item Sora
			\item Whisper
		\end{enumerate}
		
		\vspace{0.1cm}
		\textcolor{myblue}{\textbf{ANSWER: A}}
		
		\vspace{0.1cm}
		\textbf{DETAILED EXPLANATION:}
		\begin{itemize}
			\item Codex by OpenAI
			\item Code generation from descriptions
			\item Natural language to code
			\item GitHub Copilot uses Codex
		\end{itemize}
		
		\vspace{0.1cm}
		\textcolor{myred}{\textbf{WHY OTHER OPTIONS ARE WRONG:}}
		\begin{itemize}
			\item \textcolor{myred}{B:} DALL-E generates images
			\item \textcolor{myred}{C:} Sora generates videos
			\item \textcolor{myred}{D:} Whisper converts speech to text
		\end{itemize}
		
		\textcolor{myred}{\textbf{EXAM TRAP:}} Codex = \textcolor{myblue}{code generation} from OpenAI!
	\end{frame}
	
	\begin{frame}{Q90: GitHub Copilot - Tutorial}
		\fontsize{6.5}{7.5}\selectfont
		\textbf{QUESTION: What is GitHub Copilot?}
		
		\vspace{0.1cm}
		\textbf{OPTIONS:}
		\begin{enumerate}
			\item An image generation tool
			\item An AI pair programmer for code suggestions
			\item A speech recognition tool
			\item A video generation tool
		\end{enumerate}
		
		\vspace{0.1cm}
		\textcolor{myblue}{\textbf{ANSWER: B}}
		
		\vspace{0.1cm}
		\textbf{DETAILED EXPLANATION:}
		\begin{itemize}
			\item AI pair programmer
			\item Code suggestions
			\item Based on Codex/OpenAI
			\item Integrated with GitHub
		\end{itemize}
		
		\vspace{0.1cm}
		\textcolor{myred}{\textbf{WHY OTHER OPTIONS ARE WRONG:}}
		\begin{itemize}
			\item \textcolor{myred}{A:} Image generation is different
			\item \textcolor{myred}{C:} Speech recognition is different
			\item \textcolor{myred}{D:} Video generation is different
		\end{itemize}
		
		\textcolor{myred}{\textbf{EXAM TRAP:}} GitHub Copilot = \textcolor{myblue}{AI pair programmer}!
	\end{frame}
	
	\begin{frame}{Q91: Fraud Detection with GenAI - Tutorial}
		\fontsize{6.5}{7.5}\selectfont
		\textbf{QUESTION: How does GenAI help with fraud detection?}
		
		\vspace{0.1cm}
		\textbf{OPTIONS:}
		\begin{enumerate}
			\item By ignoring transactions
			\item By automating monitoring and identifying anomalies
			\item By removing all alerts
			\item By approving all transactions
		\end{enumerate}
		
		\vspace{0.1cm}
		\textcolor{myblue}{\textbf{ANSWER: B}}
		
		\vspace{0.1cm}
		\textbf{DETAILED EXPLANATION:}
		\begin{itemize}
			\item Automate monitoring
			\item Compile transaction data
			\item Identify anomalies/outliers
			\item Real-time detection
		\end{itemize}
		
		\vspace{0.1cm}
		\textcolor{myred}{\textbf{WHY OTHER OPTIONS ARE WRONG:}}
		\begin{itemize}
			\item \textcolor{myred}{A:} Ignoring transactions would miss fraud
			\item \textcolor{myred}{C:} Removing alerts misses fraud
			\item \textcolor{myred}{D:} Approving all transactions allows fraud
		\end{itemize}
		
		\textcolor{myred}{\textbf{EXAM TRAP:}} GenAI = \textcolor{myblue}{automated fraud detection}!
	\end{frame}
	
	\begin{frame}{Q92: Mastercard Decision Intelligence Pro - Tutorial}
		\fontsize{6.5}{7.5}\selectfont
		\textbf{QUESTION: What is Mastercard's Decision Intelligence Pro?}
		
		\vspace{0.1cm}
		\textbf{OPTIONS:}
		\begin{enumerate}
			\item A general-purpose chatbot
			\item A dedicated GenAI fraud detection tool
			\item A video generation tool
			\item A speech recognition tool
		\end{enumerate}
		
		\vspace{0.1cm}
		\textcolor{myblue}{\textbf{ANSWER: B}}
		
		\vspace{0.1cm}
		\textbf{DETAILED EXPLANATION:}
		\begin{itemize}
			\item Dedicated GenAI fraud detection
			\item Advanced pattern recognition
			\item Financial transaction monitoring
			\item Real-time detection
		\end{itemize}
		
		\vspace{0.1cm}
		\textcolor{myred}{\textbf{WHY OTHER OPTIONS ARE WRONG:}}
		\begin{itemize}
			\item \textcolor{myred}{A:} It's specialized, not general-purpose
			\item \textcolor{myred}{C:} It's for fraud detection, not video
			\item \textcolor{myred}{D:} It's for fraud detection, not speech
		\end{itemize}
		
		\textcolor{myred}{\textbf{EXAM TRAP:}} Decision Intelligence Pro = \textcolor{myblue}{fraud detection tool}!
	\end{frame}
	
	\begin{frame}{Q93: Cybersecurity with LLMs - Tutorial}
		\fontsize{6.5}{7.5}\selectfont
		\textbf{QUESTION: How do LLMs help with cybersecurity?}
		
		\vspace{0.1cm}
		\textbf{OPTIONS:}
		\begin{enumerate}
			\item By ignoring all threats
			\item By monitoring for potential threats and identifying vulnerabilities
			\item By removing all security measures
			\item By approving all access requests
		\end{enumerate}
		
		\vspace{0.1cm}
		\textcolor{myblue}{\textbf{ANSWER: B}}
		
		\vspace{0.1cm}
		\textbf{DETAILED EXPLANATION:}
		\begin{itemize}
			\item Monitor emails and data streams
			\item Identify potential cyber threats
			\item Cross-reference with system profile
			\item Pinpoint vulnerabilities
		\end{itemize}
		
		\vspace{0.1cm}
		\textcolor{myred}{\textbf{WHY OTHER OPTIONS ARE WRONG:}}
		\begin{itemize}
			\item \textcolor{myred}{A:} Ignoring threats is dangerous
			\item \textcolor{myred}{C:} Security measures are essential
			\item \textcolor{myred}{D:} Approving all requests allows attacks
		\end{itemize}
		
		\textcolor{myred}{\textbf{EXAM TRAP:}} LLMs = \textcolor{myblue}{cybersecurity monitoring and threat detection}!
	\end{frame}
	
	\begin{frame}{Q94: Cisco Systems LLM - Tutorial}
		\fontsize{6.5}{7.5}\selectfont
		\textbf{QUESTION: What does Cisco Systems use LLMs for?}
		
		\vspace{0.1cm}
		\textbf{OPTIONS:}
		\begin{enumerate}
			\item Image generation
			\item Malware detection
			\item Speech recognition
			\item Video creation
		\end{enumerate}
		
		\vspace{0.1cm}
		\textcolor{myblue}{\textbf{ANSWER: B}}
		
		\vspace{0.1cm}
		\textbf{DETAILED EXPLANATION:}
		\begin{itemize}
			\item Custom LLM for malware detection
			\item Security analysis
			\item Threat identification
			\item Cybersecurity protection
		\end{itemize}
		
		\vspace{0.1cm}
		\textcolor{myred}{\textbf{WHY OTHER OPTIONS ARE WRONG:}}
		\begin{itemize}
			\item \textcolor{myred}{A:} Cisco focuses on cybersecurity, not images
			\item \textcolor{myred}{C:} Cisco focuses on malware, not speech
			\item \textcolor{myred}{D:} Cisco focuses on detection, not video
		\end{itemize}
		
		\textcolor{myred}{\textbf{EXAM TRAP:}} Cisco = \textcolor{myblue}{malware detection} with LLMs!
	\end{frame}
	
	\begin{frame}{Q95: CrowdStrike Charlotte AI - Tutorial}
		\fontsize{6.5}{7.5}\selectfont
		\textbf{QUESTION: What is CrowdStrike's Charlotte AI used for?}
		
		\vspace{0.1cm}
		\textbf{OPTIONS:}
		\begin{enumerate}
			\item Image generation
			\item Cybersecurity and threat response
			\item Speech recognition
			\item Video creation
		\end{enumerate}
		
		\vspace{0.1cm}
		\textcolor{myblue}{\textbf{ANSWER: B}}
		
		\vspace{0.1cm}
		\textbf{DETAILED EXPLANATION:}
		\begin{itemize}
			\item Charlotte AI for cybersecurity
			\item Automated threat response
			\item Security operations
			\item Pattern recognition
		\end{itemize}
		
		\vspace{0.1cm}
		\textcolor{myred}{\textbf{WHY OTHER OPTIONS ARE WRONG:}}
		\begin{itemize}
			\item \textcolor{myred}{A:} Charlotte AI is for security, not images
			\item \textcolor{myred}{C:} Charlotte AI is for threat detection, not speech
			\item \textcolor{myred}{D:} Charlotte AI is for security, not video
		\end{itemize}
		
		\textcolor{myred}{\textbf{EXAM TRAP:}} Charlotte AI = \textcolor{myblue}{cybersecurity threat response}!
	\end{frame}
	
	\begin{frame}{Q96: EXL Insurance LLM - Tutorial}
		\fontsize{6.5}{7.5}\selectfont
		\textbf{QUESTION: What is the EXL Insurance LLM used for?}
		
		\vspace{0.1cm}
		\textbf{OPTIONS:}
		\begin{enumerate}
			\item General purpose chat
			\item Claims processing, underwriting, and customer service
			\item Image generation
			\item Video creation
		\end{enumerate}
		
		\vspace{0.1cm}
		\textcolor{myblue}{\textbf{ANSWER: B}}
		
		\vspace{0.1cm}
		\textbf{DETAILED EXPLANATION:}
		\begin{itemize}
			\item Claims processing
			\item Underwriting
			\item Customer service
			\item Insurance-specific tasks
		\end{itemize}
		
		\vspace{0.1cm}
		\textcolor{myred}{\textbf{WHY OTHER OPTIONS ARE WRONG:}}
		\begin{itemize}
			\item \textcolor{myred}{A:} EXL is insurance-specific, not general
			\item \textcolor{myred}{C:} EXL is text-based, not image
			\item \textcolor{myred}{D:} EXL is document processing, not video
		\end{itemize}
		
		\textcolor{myred}{\textbf{EXAM TRAP:}} EXL Insurance LLM = \textcolor{myblue}{claims + underwriting + customer service}!
	\end{frame}
	
	\begin{frame}{Q97: EXL Insurance LLM Results - Tutorial}
		\fontsize{6.5}{7.5}\selectfont
		\textbf{QUESTION: What performance improvement did EXL Insurance LLM achieve?}
		
		\vspace{0.1cm}
		\textbf{OPTIONS:}
		\begin{enumerate}
			\item 10\% reduction in operational costs
			\item 30\% reduction in operational costs
			\item 50\% reduction in operational costs
			\item No reduction in operational costs
		\end{enumerate}
		
		\vspace{0.1cm}
		\textcolor{myblue}{\textbf{ANSWER: B}}
		
		\vspace{0.1cm}
		\textbf{DETAILED EXPLANATION:}
		\begin{itemize}
			\item 30\% reduction in operational costs
			\item Through efficient claims processing
			\item Streamlined underwriting
			\item Automated customer service
		\end{itemize}
		
		\vspace{0.1cm}
		\textcolor{myred}{\textbf{WHY OTHER OPTIONS ARE WRONG:}}
		\begin{itemize}
			\item \textcolor{myred}{A:} Too low compared to actual results
			\item \textcolor{myred}{C:} Too high compared to actual results
			\item \textcolor{myred}{D:} There was a significant reduction
		\end{itemize}
		
		\textcolor{myred}{\textbf{EXAM TRAP:}} EXL Insurance LLM = \textcolor{myblue}{30\% cost reduction}!
	\end{frame}
	
	\begin{frame}{Q98: Investment Firms Using LLMs - Tutorial}
		\fontsize{6.5}{7.5}\selectfont
		\textbf{QUESTION: Which investment firms are using LLMs?}
		
		\vspace{0.1cm}
		\textbf{OPTIONS:}
		\begin{enumerate}
			\item Only BlackRock
			\item BlackRock, Bridgewater, and AQR
			\item Only Bridgewater
			\item No investment firms use LLMs
		\end{enumerate}
		
		\vspace{0.1cm}
		\textcolor{myblue}{\textbf{ANSWER: B}}
		
		\vspace{0.1cm}
		\textbf{DETAILED EXPLANATION:}
		\begin{itemize}
			\item \textbf{BlackRock:} LLMs in investment process
			\item \textbf{Bridgewater:} \$2 billion ML fund
			\item \textbf{AQR:} LLM integration for strategies
		\end{itemize}
		
		\vspace{0.1cm}
		\textcolor{myred}{\textbf{WHY OTHER OPTIONS ARE WRONG:}}
		\begin{itemize}
			\item \textcolor{myred}{A:} More firms than just BlackRock
			\item \textcolor{myred}{C:} More firms than just Bridgewater
			\item \textcolor{myred}{D:} Many firms use LLMs
		\end{itemize}
		
		\textcolor{myred}{\textbf{EXAM TRAP:}} BlackRock, Bridgewater, AQR = \textcolor{myblue}{all using LLMs}!
	\end{frame}
	
	\begin{frame}{Q99: GenAI Impact Summary - Tutorial}
		\fontsize{6.5}{7.5}\selectfont
		\textbf{QUESTION: What areas of finance are being transformed by GenAI?}
		
		\vspace{0.1cm}
		\textbf{OPTIONS:}
		\begin{enumerate}
			\item Only fraud detection
			\item Fraud detection, customer service, document processing, and investment
			\item Only customer service
			\item Only document processing
		\end{enumerate}
		
		\vspace{0.1cm}
		\textcolor{myblue}{\textbf{ANSWER: B}}
		
		\vspace{0.1cm}
		\textbf{DETAILED EXPLANATION:}
		\begin{itemize}
			\item \textbf{Fraud Detection:} Automated monitoring
			\item \textbf{Customer Service:} Chatbots
			\item \textbf{Document Processing:} Summarization
			\item \textbf{Investment:} Market analysis and trading
		\end{itemize}
		
		\vspace{0.1cm}
		\textcolor{myred}{\textbf{WHY OTHER OPTIONS ARE WRONG:}}
		\begin{itemize}
			\item \textcolor{myred}{A:} GenAI impacts multiple areas
			\item \textcolor{myred}{C:} GenAI impacts more than customer service
			\item \textcolor{myred}{D:} GenAI impacts more than document processing
		\end{itemize}
		
		\textcolor{myred}{\textbf{EXAM TRAP:}} GenAI transforms \textcolor{myblue}{multiple financial sectors}!
	\end{frame}
	
	\begin{frame}{Q100: Final Summary - Tutorial}
		\fontsize{6.5}{7.5}\selectfont
		\textbf{QUESTION: What is the most comprehensive summary of Generative AI?}
		
		\vspace{0.1cm}
		\textbf{OPTIONS:}
		\begin{enumerate}
			\item GenAI is only about text generation
			\item GenAI creates new content, uses Word2Vec/RNNs/Transformers/LLMs, requires careful evaluation, and has broad applications
			\item GenAI is only about image generation
			\item GenAI doesn't require evaluation
		\end{enumerate}
		
		\vspace{0.1cm}
		\textcolor{myblue}{\textbf{ANSWER: B}}
		
		\vspace{0.1cm}
		\textbf{DETAILED EXPLANATION:}
		\begin{itemize}
			\item \textbf{Creates new content:} Text, images, audio, video
			\item \textbf{Key technologies:} Word2Vec, RNNs, Transformers, LLMs
			\item \textbf{Evaluation:} Quality, hallucination, robustness, fairness
			\item \textbf{Applications:} Broad across all financial sectors
		\end{itemize}
		
		\vspace{0.1cm}
		\textcolor{myred}{\textbf{WHY OTHER OPTIONS ARE WRONG:}}
		\begin{itemize}
			\item \textcolor{myred}{A:} GenAI includes more than text
			\item \textcolor{myred}{C:} GenAI includes more than images
			\item \textcolor{myred}{D:} Evaluation is essential for GenAI
		\end{itemize}
		
		\textcolor{myred}{\textbf{EXAM SUCCESS:}} Understand \textcolor{myblue}{technology, evaluation, and applications}!
	\end{frame}
	
\end{document}